\documentclass[12pt,oneside]{report}

\usepackage[spanish,es-nodecimaldot]{babel}
\usepackage[utf8]{inputenc}
\usepackage[T1]{fontenc}
\usepackage[a4paper,margin=2.5cm]{geometry}

\usepackage{csquotes}
\usepackage{graphicx}
\usepackage{booktabs}
\usepackage{hyperref}
\usepackage{amsmath, amssymb}
\usepackage{caption}
\usepackage{listings}
\usepackage{algorithm}
\usepackage{algpseudocode}
\usepackage{cite}
\usepackage{array}
\usepackage{colortbl}
\usepackage{xcolor}
\usepackage{float}
\usepackage{tikz}
\usepackage{longtable}
\usepackage{needspace}
\usepackage{multirow}
\usepackage{tabularx}
\usepackage{xltabular}
\usetikzlibrary{arrows.meta,positioning,shapes.geometric}
\urlstyle{same}

% ============================================
% CORRECCIONES PARA EVITAR ESPACIOS EN BLANCO
% ============================================
\raggedbottom  % No estira verticalmente el contenido de cada página

% Parámetros de floats más permisivos
\renewcommand{\topfraction}{0.9}       % hasta 90% de la página puede ser floats arriba
\renewcommand{\bottomfraction}{0.9}    % hasta 90% abajo
\renewcommand{\textfraction}{0.07}     % mínimo 7% de texto en una página con floats
\renewcommand{\floatpagefraction}{0.7} % una página de solo floats necesita al menos 70%
\setcounter{topnumber}{4}              % hasta 4 floats arriba por página
\setcounter{bottomnumber}{4}           % hasta 4 floats abajo por página
\setcounter{totalnumber}{8}            % hasta 8 floats totales por página

% Reducir penalidades para que LaTeX no deje huecos
\setlength{\intextsep}{8pt plus 2pt minus 2pt}    % espacio alrededor de floats [H]
\setlength{\floatsep}{8pt plus 2pt minus 2pt}      % entre floats consecutivos
\setlength{\textfloatsep}{10pt plus 2pt minus 4pt} % entre float y texto

% Definición de colores para código
\definecolor{lightgray}{rgb}{.9,.9,.9}
\definecolor{darkgray}{rgb}{.4,.4,.4}
\definecolor{purple}{rgb}{0.65, 0.12, 0.82}

% Colores para matrices de correlación
\definecolor{expectedGreen}{HTML}{D4EDDA}
\definecolor{suspiciousYellow}{HTML}{FFF3CD}
\definecolor{incoherentRed}{HTML}{F8D7DA}
\definecolor{headerBlue}{HTML}{2E75B6}

% Configuración de Listings para JavaScript
\lstdefinelanguage{JavaScript}{
  keywords={typeof, new, true, false, catch, function, return, null, catch, switch, var, if, in, while, do, else, case, break},
  keywordstyle=\color{blue}\bfseries,
  ndkeywords={class, export, boolean, throw, implements, import, this},
  ndkeywordstyle=\color{darkgray}\bfseries,
  identifierstyle=\color{black},
  sensitive=false,
  comment=[l]{//},
  morecomment=[s]{/*}{*/},
  commentstyle=\color{purple}\ttfamily,
  stringstyle=\color{red}\ttfamily,
  morestring=[b]',
  morestring=[b]"
}

\lstset{
  language=JavaScript,
  basicstyle=\ttfamily\footnotesize,
  keywordstyle=\color{blue},
  commentstyle=\color{green!60!black},
  stringstyle=\color{red},
  numbers=left,
  numberstyle=\tiny\color{gray},
  stepnumber=1,
  numbersep=5pt,
  backgroundcolor=\color{white},
  showspaces=false,
  showstringspaces=false,
  showtabs=false,
  frame=single,
  tabsize=2,
  captionpos=b,
  breaklines=true,
  breakatwhitespace=false,
  escapeinside={\%*}{*)}
}

\hypersetup{colorlinks=true,linkcolor=blue,citecolor=blue,urlcolor=blue}

\begin{document}

% =========================
% PORTADA
% =========================
\begin{titlepage}
\newgeometry{top=3.0cm,bottom=3.2cm,left=3.0cm,right=3.0cm}
\centering
\setlength{\parindent}{0pt}

\vspace{0.9cm}
{\Large \textbf{Universidad de los Andes}\par}
{\large Departamento de Ingeniería de Sistemas y Computación\par}

\vspace{2.6cm}
{\bfseries\LARGE Control de Acceso Basado en Contexto para Extensiones de Navegador Web\par}
\vspace{0.45cm}
{\large \textit{Proyecto de Grado}\par}

\vspace{9.6cm}
\begin{flushleft}
\renewcommand{\arraystretch}{1.35}
\begin{tabular}{@{}l l@{}}
Autor: & Adriana Sofia Rozo Cepeda \\
Asesor: & Yuri Andrea Pinto Rojas\\
\end{tabular}
\end{flushleft}

\vfill
\begin{flushleft}
Bogotá, Colombia\\
\today
\end{flushleft}

\restoregeometry
\end{titlepage}
% --- TABLA DE CONTENIDO ---
\tableofcontents
\clearpage % Asegura que la introducción empiece en una página nueva


% ============================================
% CAPÍTULO 1: INTRODUCCIÓN
% ============================================

\chapter{Introducción}

% =========================
% PROBLEMA
% =========================
\section{Planteamiento del Problema}
\label{sec:problema}

\begingroup
\setlength{\parindent}{0pt}
\setlength{\parskip}{0.6em}

En la actualidad, el navegador web se ha consolidado como la plataforma central de la vida digital, donde se procesan desde transacciones bancarias y datos corporativos confidenciales hasta interacciones con modelos de inteligencia artificial. Uno de sus componentes más usados son las extensiones que fueron creadas para personalizar la experiencia de usuario y extender las funcionalidades del navegador; sin embargo, dicha versatilidad ha transformado a las mismas en uno de los vectores de ataque más críticos y menos gestionados del ecosistema de ciberseguridad actual. El problema reside en que las extensiones operan dentro del contexto de confianza del usuario, pero su modelo de seguridad presenta tres deficiencias:

\begin{enumerate}
    \item Opacidad en el modelo de privilegios: El ecosistema actual de navegadores basados en Chromium permite que las extensiones soliciten permisos excesivamente amplios que el usuario promedio acepta sin comprender las implicaciones reales. No existen mecanismos amigables con el usuario que evalúen de forma automática si la funcionalidad declarada de una extensión es coherente con los permisos técnicos que solicita.
    \item Actualizaciones sin verificación de seguridad: El navegador hace efectivas las actualizaciones de forma silenciosa y automática y los cibercriminales aprovechan esto para introducir malware.
    \item Ausencia de un control de acceso contextual: Las extensiones operan bajo un esquema de ejecución permanente de todo o nada. Esto implica que una extensión diseñada para el entretenimiento o la productividad, mantiene el mismo nivel de privilegios cuando el usuario navega en sitios recreativos que cuando accede a entornos de alta criticidad, como portales bancarios, sistemas de salud o herramientas corporativas. Esta falta de aislamiento facilita la exfiltración de información sensible por parte de extensiones que no deberían tener actividad en estos contextos críticos.
\end{enumerate}

En consecuencia, existe una brecha de seguridad significativa ya que los usuarios carecen de herramientas integradas en el navegador que supervisen el comportamiento de las extensiones en tiempo real según el contexto de navegación y la evolución de sus permisos.

\endgroup

% =========================
% JUSTIFICACIÓN
% =========================
\section{Justificación}
\label{sec:justificacion}

\begingroup
\setlength{\parindent}{0pt}
\setlength{\parskip}{0.6em}

Este proyecto de grado busca cerrar la brecha entre el desconocimiento de los riesgos y las vulnerabilidades de las extensiones por parte del usuario y la sofisticación de los ataques actuales.

En primer lugar, desde una perspectiva técnica, la transición a Manifest V3, aunque positiva en cuanto a la eliminación de código remoto arbitrario, no ha resuelto el problema de la confianza ciega. Como evidencian incidentes recientes, los mecanismos de revisión de las tiendas oficiales son incapaces de detectar comportamiento malicioso introducido posterior a la instalación \cite{thn2025shadypanda}. Más aún, investigaciones sobre el abuso de APIs nativas en Manifest V3 han demostrado que los atacantes ahora utilizan funcionalidades legítimas, como \texttt{chrome.cookies} para exfiltración e inyección de tokens de sesión en campañas como DataByCloud \cite{pandya2026databycloud}, o \texttt{chrome.scripting} para inyectar keyloggers empaquetados en lugar de descargar código remoto \cite{singh2025study}.

Ahora bien, las herramientas de seguridad tradicionales (antivirus de red, firewalls) no tienen visibilidad dentro del proceso aislado del navegador debido al protocolo HTTPS y la arquitectura de \textit{sandboxing} de las extensiones. Por tanto, es necesario desarrollar una solución que opere desde el interior del propio navegador para monitorear y restringir privilegios, utilizando las APIs nativas de gestión sin requerir software externo o privilegios elevados del sistema operativo.

Finalmente, el impacto social y corporativo es crítico. Por ejemplo, en reportes recientes se documenta que extensiones de ``privacidad'' como Urban VPN Proxy interceptaron y vendieron conversaciones completas con herramientas de IA (ChatGPT, Claude, Gemini, entre otras), comprometiendo propiedad intelectual y datos personales de millones de usuarios \cite{dardikman2025aiconversations}. Adicionalmente, investigadores de seguridad han demostrado que las extensiones de navegador con capacidades agénticas, como los nuevos asistentes de IA integrados en el navegador, introducen superficies de ataque novedosas que incluyen inyección indirecta de prompts y exposición de tokens de sesión \cite{zenitylabs2025claude}. Este proyecto devuelve el control al usuario final con una herramienta de ``defensa activa'', permitiéndole aplicar políticas de aislamiento contextual sin necesidad de conocimientos avanzados de ciberseguridad.

\endgroup

% =========================
% OBJETIVOS
% =========================
\section{Objetivos}
\label{sec:objetivos}
\begingroup
\setlength{\parindent}{0pt}
\setlength{\parskip}{0.6em}

\subsection{Objetivo General}
Desarrollar un mecanismo de seguridad para extensiones de navegador web, que sensibilice y devuelva el control al usuario final, permitiéndole monitorear y restringir los privilegios de las extensiones según el contexto de navegación para mitigar riesgos de privacidad.

\subsection{Objetivos Específicos}
\begin{itemize}
    \item Definir un modelo de evaluación de riesgo para extensiones bajo la arquitectura Manifest V3, que identifique incoherencias entre los permisos solicitados y la funcionalidad declarada para reducir la opacidad del modelo de privilegios.
    \item Diseñar e implementar un motor de monitoreo utilizando APIs nativas capaz de detectar cambios críticos en el comportamiento o en los privilegios de las extensiones tras actualizaciones automáticas.
    \item Construir un sistema de control de acceso contextual que habilite al usuario para configurar políticas de aislamiento, restringiendo la ejecución de extensiones en entornos críticos (banca, salud, corporativo) frente a entornos de navegación general.
    \item Validar la eficacia del mecanismo propuesto mediante escenarios de prueba controlados que simulen inyección de código y exfiltración de datos, evaluando tanto la capacidad de detección técnica como la usabilidad de las alertas para el usuario.
\end{itemize}
\endgroup


% ============================================
% CAPÍTULO 2: MARCO TEÓRICO
% ============================================

\chapter{Marco Teórico}
\begingroup
\setlength{\parindent}{0pt}
\setlength{\parskip}{0.6em}

\section{Contexto Técnico: Arquitectura de Extensiones de Navegador}

\subsection{Modelo de Extensiones en Chromium}

Las extensiones de navegador en Google Chrome son aplicaciones web empaquetadas que extienden la funcionalidad del navegador mediante JavaScript, HTML y CSS \cite{chrome_extensions_overview}. Si bien otros navegadores basados en Chromium (Microsoft Edge, Brave, Opera) comparten la misma arquitectura de extensiones, este proyecto se centra exclusivamente en Google Chrome, dado que es el navegador con mayor cuota de mercado, su Chrome Web Store es la plataforma de distribución principal, y la validación experimental se realizará únicamente en este navegador. La arquitectura actual se basa en un modelo de aislamiento que separa scripts según su nivel de privilegio \cite{chrome_extensions_security}:

\begin{itemize}
    \item \textbf{Service Workers} (anteriormente ``background scripts'' en MV2): Scripts con acceso completo a las Chrome Extension APIs privilegiadas (cookies, historial, pestañas). Operan en un contexto aislado sin acceso directo al DOM de las páginas web.
    \item \textbf{Content Scripts}: Scripts que se inyectan en páginas web y tienen acceso al DOM, pero con acceso limitado a Extension APIs. Pueden comunicarse con Service Workers mediante \textit{message passing}.
    \item \textbf{Scripts de página web}: Scripts ejecutados por sitios web, completamente aislados de la extensión por el principio de \textit{isolated worlds} \cite{isolated_worlds}.
\end{itemize}

La comunicación entre estos componentes se realiza mediante canales de mensajería (\texttt{chrome.runtime.sendMessage}, \texttt{chrome.runtime.onMessage}) que constituyen puntos críticos de seguridad, ya que datos no sanitizados provenientes de páginas web pueden fluir hacia APIs privilegiadas.

\subsection{Transición de Manifest V2 a Manifest V3}

Manifest V3 fue lanzado en Chrome 88 (enero de 2021) y la deprecación de Manifest V2 comenzó formalmente en 2023, representando un cambio significativo más no suficiente en la arquitectura de seguridad de extensiones \cite{pantelaios2024manifestv3, manifestv3chrome}. Los cambios principales incluyen:

\begin{enumerate}
    \item \textbf{Prohibición de código remoto}: Las extensiones ya no pueden ejecutar código descargado dinámicamente mediante \texttt{eval()}, \texttt{new Function()}, o scripts externos no declarados en el manifest \cite{manifestv3chrome}.
    \item \textbf{Restricción de webRequest}: La capacidad de bloqueo de la API \texttt{webRequest} (\texttt{webRequestBlocking}) fue restringida en MV3 exclusivamente a extensiones instaladas por política empresarial. Para extensiones regulares, fue reemplazada por \texttt{declarativeNetRequest}, que permite modificar peticiones mediante reglas declarativas sin exponer el contenido completo del tráfico \cite{declarative_net_request}.
    \item \textbf{Service Workers con ciclo de vida limitado}: Los background scripts persistentes fueron reemplazados por Service Workers que se activan solo cuando es necesario y se suspenden tras inactividad, reduciendo el consumo de recursos \cite{service_workers}. Si bien desde Chrome 110 existen mecanismos que permiten extender la vida del Service Worker, como \texttt{chrome.runtime.connectNative()}, que mantiene el worker activo mientras exista una conexión nativa, estos requieren la instalación de un binario en el sistema operativo del usuario, lo que eleva significativamente la barrera de entrada para atacantes. En la práctica, el vector dominante de persistencia en extensiones maliciosas es \texttt{chrome.alarms}, que permite crear ``heartbeats'' periódicos sin requerir componentes externos \cite{seetharam2025genai, socket2025phantom}.
\end{enumerate}

Sin embargo, investigaciones recientes demuestran que MV3 no ha eliminado las vulnerabilidades de extensiones; los atacantes simplemente han adaptado sus técnicas para abusar de las APIs permitidas.

\section{Vectores de Ataque Documentados}

\subsection{Sleeper Extensions y Ataques de Cadena de Suministro}

El vector de ataque más sofisticado documentado en 2024-2025 es el uso de \textit{sleeper extensions}: extensiones legítimas que operan sin comportamiento malicioso durante períodos prolongados (meses o años) para acumular usuarios y reputación, antes de ser ``militarizadas'' mediante una actualización que introduce código malicioso.

\textbf{Caso de estudio: DarkSpectre}

La operación DarkSpectre, desenmascarada a finales de diciembre de 2025 por Koi Security, comprometió a 8.8 millones de usuarios mediante tres campañas coordinadas: ShadyPanda (5.6 millones de usuarios), GhostPoster (1.05 millones) y The Zoom Stealer (2.2 millones) \cite{koi2025darkspectre}. El modus operandi consistía en:

\begin{enumerate}
    \item Publicar extensiones funcionales (calculadoras, convertidores de unidades, ad blockers) en Chrome Web Store.
    \item Operar legítimamente durante 2-4 años, acumulando usuarios y reviews positivas.
    \item Adquirir o comprometer la cuenta del desarrollador original.
    \item Publicar una actualización que introduce código de exfiltración, el cual se ejecuta automáticamente en todos los dispositivos sin intervención del usuario.
\end{enumerate}

Este patrón aprovecha una vulnerabilidad crítica: las actualizaciones automáticas de Chrome se aplican sin intervención del usuario siempre que no se añadan permisos que generen un \textit{warning} de nivel superior al ya concedido. En la práctica, los atacantes operan dentro de los permisos ya otorgados, modificando el comportamiento del código sin alterar el manifest, lo que les permite introducir funcionalidad maliciosa sin que el usuario sea notificado. Pantelaios et al. documentaron sistemáticamente este fenómeno de ``cambio malicioso'' post-publicación, demostrando que los deltas de actualización son un indicador eficaz para detectar extensiones comprometidas \cite{pantelaios2020malicious}. Más recientemente, en 2026, se reportó que extensiones de estilo visual infectaron a 500,000 usuarios de VKontakte mediante este mismo vector \cite{kathir2026vkstyles}.

\subsection{Abuso de APIs de Manifest V3}

Contrario a la expectativa de que MV3 eliminaría las extensiones maliciosas, reportes de 2025 documentan el abuso de APIs legítimas \cite{seetharam2025genai}:

\textbf{Abuso de chrome.cookies y chrome.scripting (Campaña DataByCloud):}

La campaña DataByCloud, documentada por Socket Threat Research en enero de 2026, desplegó cinco extensiones coordinadas que utilizaban las APIs \texttt{chrome.cookies} y \texttt{chrome.scripting} para exfiltrar tokens de autenticación de plataformas corporativas como Workday, NetSuite y SuccessFactors \cite{pandya2026databycloud}. Las extensiones solicitaban permisos de \texttt{cookies}, \texttt{management}, \texttt{scripting}, \texttt{storage} y \texttt{declarativeNetRequest} con \texttt{host\_permissions} para los dominios objetivo. El mecanismo de ataque consistía en:

\begin{enumerate}
    \item \textbf{Exfiltración continua:} Las extensiones extraían cookies de autenticación (\texttt{\_\_session}) de los dominios objetivo mediante \texttt{chrome.cookies.get()} y las transmitían al servidor de comando y control cada 60 segundos.
    \item \textbf{Inyección de sesiones:} La extensión más sofisticada (Software Access) podía recibir cookies robadas desde el servidor del atacante e inyectarlas en el navegador de un tercero mediante \texttt{chrome.cookies.set()}, permitiendo el secuestro directo de sesiones sin necesidad de credenciales o autenticación multifactor.
    \item \textbf{Bloqueo de respuesta a incidentes:} Dos extensiones manipulaban el DOM para bloquear el acceso a 44-56 páginas administrativas de Workday (cambio de contraseñas, gestión de 2FA, políticas de seguridad), impidiendo que los equipos de seguridad remediaran el acceso no autorizado.
\end{enumerate}

Este ataque no requería \texttt{eval()} ni código remoto, y todas las APIs utilizadas son legítimas en MV3, pero su combinación coordinada con \texttt{host\_permissions} amplios permitía el compromiso completo de cuentas empresariales.

\textbf{Persistencia mediante chrome.alarms y otros mecanismos:}

Dado que los Service Workers de MV3 tienen ciclo de vida limitado, los atacantes emplean múltiples estrategias para mantener persistencia:
\begin{itemize}
    \item \textbf{\texttt{chrome.alarms}}: El vector dominante documentado en campañas reales. Crea ``heartbeats'' que despiertan al Service Worker periódicamente sin dependencias externas \cite{seetharam2025genai, stealth_extension}.
    \item \textbf{\texttt{chrome.runtime.connect} desde content scripts}: Mantener un puerto de comunicación abierto entre un content script y el Service Worker impide que Chrome lo suspenda. Requiere \texttt{host\_permissions} y al menos una pestaña activa en un dominio con permiso.
    \item \textbf{\texttt{chrome.offscreen}}: Permite crear un documento HTML oculto con acceso completo al DOM, donde se pueden ejecutar \texttt{setInterval}, WebSockets persistentes y otras APIs que el Service Worker no soporta. Chrome limita a un documento offscreen por extensión.
\end{itemize}

\subsection{Exfiltración de Conversaciones con Herramientas de IA}

Un vector emergente es la interceptación de conversaciones con plataformas de inteligencia artificial como ChatGPT, Claude y Gemini \cite{dardikman2025aiconversations}. Extensiones que se anuncian como ``mejoras de privacidad'' o ``VPNs gratuitos'' inyectan content scripts que:

\begin{enumerate}
    \item Escuchan eventos de envío de prompts mediante \texttt{addEventListener('submit')}.
    \item Capturan el contenido del DOM donde se muestran las respuestas de la IA.
    \item Exfiltran ambos (prompt y respuesta) a servidores externos, comprometiendo propiedad intelectual, secretos comerciales, y datos personales sensibles.
\end{enumerate}

Koi Research reporta que esta práctica comprometió conversaciones de 8 millones de usuarios \cite{dardikman2025aiconversations}.

\section{Trabajos Relacionados}

\subsection{Análisis Estático de Extensiones}

\textbf{CoCo: Concurrent Abstract Interpretation}

CoCo es una herramienta de investigación que utiliza interpretación abstracta concurrente para detectar vulnerabilidades en extensiones de Chrome. Sin embargo, CoCo opera fuera del navegador, requiere el código fuente completo de la extensión, y está diseñado para análisis en batch en entornos de investigación, no para protección de un usuario final no académico \cite{yu2023coco}.

\textbf{DoubleX: Extension Dependence Graph}

DoubleX mejora sobre trabajos previos (EmPoWeb) mediante un grafo de dependencias de extensión (EDG) para automatizar la detección de vulnerabilidades. Sin embargo, no puede manejar características dinámicas de JavaScript como llamadas a funciones resueltas en tiempo de ejecución \cite{doublex}.

\textbf{Arcanum: Privacy Risk Detection}

Arcanum detecta riesgos de privacidad mediante análisis de flujo de información dinámico (dynamic taint analysis). Rastrea datos sensibles desde su origen (DOM, cookies) hasta potenciales puntos de exfiltración (peticiones de red). Su principal limitación es que requiere instrumentación del motor Chromium, haciéndolo inviable para uso comercial \cite{arcanum}.

\subsection{Sistemas de Monitoreo en Tiempo Real}

\textbf{Mitigating Security Risks (Sam et al.) \cite{sam2023mitigating}}

Sam y Jenifer proponen una extensión de monitoreo que analiza otras extensiones en tiempo real. Su arquitectura incluye análisis de vulnerabilidades mediante patrones conocidos, monitoreo de comportamiento en runtime, sistema de educación al usuario. 

Sin embargo, el trabajo presenta la arquitectura a nivel conceptual, pero no especifica los algoritmos concretos de detección ni los modelos utilizados. Tampoco implementa mecanismos preventivos como cuarentena de actualizaciones. Además, la evaluación experimental carece de métricas cuantitativas rigurosas (precisión, recall, tasa de falsos positivos), lo que limita la validación empírica del enfoque.

\subsection{Tabla Comparativa de Trabajos Relacionados}

\begin{table}[!htbp]
\centering
\small
\begin{tabular}{|p{2.5cm}|p{2.5cm}|p{3cm}|p{3cm}|p{3cm}|}
\hline
\textbf{Trabajo} & \textbf{Enfoque} & \textbf{Fortalezas} & \textbf{Limitaciones} & \textbf{Diferencia con esta propuesta} \\
\hline
CoCo \cite{yu2023coco} & Interpretación abstracta concurrente & Detecta zero-days, maneja JS dinámico, baja tasa FP & Opera fuera del navegador, para investigadores & Opera dentro del navegador, para usuarios finales \\
\hline
DoubleX \cite{doublex} & Extension Dependence Graph & Automatizado, mejora sobre EmPoWeb & No maneja código dinámico, alta tasa FP & Enfoque preventivo contextual vs. detección offline \\
\hline
Arcanum \cite{arcanum} & Taint analysis dinámico & Rastrea flujos sensibles, detecta exfiltración & Requiere modificar Chromium, sobrecarga & Usa APIs estándar, sin modificar navegador \\
\hline
Mitigating \cite{sam2023mitigating} & Monitoreo en tiempo real & Accesible a usuarios, opera en navegador & Algoritmos no especificados, sin cuarentena & Análisis contextual + comparación de manifests \\
\hline
\end{tabular}
\caption{Comparación de trabajos relacionados}
\end{table}

\section{Marco Regulatorio y Políticas de Plataforma}

A diferencia de otros componentes de software, las extensiones de navegador operan en un vacío regulatorio específico. No existe legislación dedicada que regule su desarrollo, distribución o comportamiento. Sin embargo, están sujetas a múltiples capas de gobernanza que se superponen parcialmente.

\subsection{Políticas de Plataforma: Chrome Web Store}

Google establece requisitos obligatorios para la publicación de extensiones mediante las \textit{Chrome Web Store Program Policies} \cite{cws_policies}. Estas políticas, actualizadas en mayo de 2025, constituyen el marco normativo más específico para extensiones y establecen:

\textbf{Requisitos de privacidad:}
\begin{itemize}
    \item Obligación de publicar una política de privacidad si la extensión maneja datos de usuario.
    \item Divulgación prominente del uso de datos antes de la instalación.
    \item Transmisión de datos sensibles exclusivamente mediante HTTPS.
    \item Prohibición de transferir datos de usuario para publicidad personalizada o a brokers de datos.
\end{itemize}

\textbf{Requisitos técnicos (Manifest V3):}
\begin{itemize}
    \item Prohibición de código ofuscado o funcionalidad oculta.
    \item El código debe ser ``fácilmente discernible'' durante la revisión.
    \item Prohibición de ejecutar lógica desde fuentes remotas (excepto APIs específicas como Debugger y User Scripts).
    \item Verificación en dos pasos obligatoria para cuentas de desarrollador.
\end{itemize}

\textbf{Limitaciones:} Estas políticas son contractuales, no legales. Google puede remover extensiones, pero no puede imponer sanciones legales. Además, las extensiones instaladas fuera del Web Store (\textit{sideloading}) no están sujetas a estas políticas.

\subsection{Regulaciones de Privacidad Aplicables}

Las extensiones que procesan datos personales de usuarios en ciertas jurisdicciones están sujetas a regulaciones generales de privacidad:

\begin{table}[!htbp]
\centering
\small
\begin{tabularx}{\textwidth}{|l|l|X|}
\hline
\textbf{Regulación} & \textbf{Jurisdicción} & \textbf{Aplicación a Extensiones} \\
\hline
GDPR (2018) & Unión Europea & Aplica si la extensión recolecta datos personales de usuarios en la UE, independientemente de dónde esté el desarrollador. Requiere base legal para el procesamiento, consentimiento informado, y derecho de acceso/eliminación. \\
\hline
CCPA/CPRA & California, EE.UU. & Aplica a extensiones que recolectan datos de residentes de California. Requiere divulgación de categorías de datos recolectados y derecho a opt-out de venta de datos. \\
\hline
ePrivacy Directive & Unión Europea & Aplica si la extensión almacena o accede a información en el dispositivo del usuario (cookies, almacenamiento local). Requiere consentimiento previo. \\
\hline
\end{tabularx}
\caption{Regulaciones de privacidad aplicables a extensiones de navegador}
\label{tab:privacy_regulations}
\end{table}

\subsection{Digital Services Act (DSA) y Digital Markets Act (DMA)}

El \textit{Digital Services Act} (DSA), en vigor desde febrero de 2024, regula plataformas de intermediación digital en la UE \cite{eu_dsa}. Si bien no regula directamente a las extensiones, impone obligaciones a la Chrome Web Store como plataforma:

\begin{itemize}
    \item Mecanismos de reporte de contenido ilegal.
    \item Transparencia en sistemas de moderación y recomendación.
    \item Prohibición de publicidad basada en perfilado con datos sensibles (categorías especiales del GDPR).
\end{itemize}

El \textit{Digital Markets Act} (DMA) designa a Google como ``gatekeeper'' y podría afectar futuras políticas de distribución de extensiones, aunque hasta marzo de 2026 no ha habido cambios directos en este ámbito.

\subsection{Propuesta Digital Omnibus (Noviembre 2025)}

La Comisión Europea propuso en noviembre de 2025 el paquete \textit{Digital Omnibus} que incluye reformas relevantes \cite{eu_digital_omnibus}:

\begin{itemize}
    \item \textbf{Señales de consentimiento en navegadores:} Los fabricantes de navegadores deberán incluir herramientas técnicas que permitan a los usuarios dar o rechazar consentimiento de forma automática, reduciendo la dependencia de banners de cookies.
    \item \textbf{Simplificación de GDPR:} Punto único de entrada para reportar brechas de seguridad.
\end{itemize}

\textbf{Implicación para extensiones:} Si se aprueba, las extensiones que manejen cookies o tracking deberán respetar las señales de consentimiento configuradas en el navegador.

\subsection{Vacío Regulatorio Identificado}

El análisis del marco regulatorio revela un \textbf{vacío significativo}: no existe ninguna regulación que exija una evaluación de riesgo estandarizada para extensiones de navegador antes de su publicación o durante su operación. Las políticas de Google son la única barrera de entrada, pero:

\begin{enumerate}
    \item No especifican una metodología de evaluación de riesgo cuantitativa.
    \item La clasificación de riesgo de permisos (whitepaper de 2019) no ha sido actualizada desde antes de Manifest V3.
    \item No existe un estándar público para correlacionar permisos con categorías funcionales.
\end{enumerate}

Este vacío justifica la necesidad del modelo de evaluación de riesgo propuesto en este proyecto, que adapta estándares reconocidos (NIST SP 800-30, ISO/IEC 27005) al contexto específico de extensiones de navegador.

\endgroup

% ============================================
% CAPÍTULO 3: METODOLOGÍA
% ============================================

\chapter{Metodología}
\label{chap:metodologia}
\begingroup
\setlength{\parindent}{0pt}
\setlength{\parskip}{0.6em}

El desarrollo se estructurará en cuatro fases metodológicas alineadas con los objetivos específicos planteados.

\section{Fase 1: Análisis y Definición del Modelo de Riesgo}
Esta fase corresponde al primer objetivo específico y se centra en establecer las reglas lógicas que el sistema utilizará para determinar si una extensión es peligrosa.
\begin{itemize}
    \item \textbf{Revisión documental de permisos (Manifest V3):} Se clasificará la totalidad de los permisos disponibles en la API de Chrome según su nivel de criticidad.
    \item \textbf{Definición de patrones de incoherencia:} Se crearán matrices de correlación entre Funcionalidad Declarada (categoría de la tienda) y Permisos Solicitados. Ejemplo: Una extensión de categoría Calculadora no debería solicitar acceso a \texttt{webRequest} o \texttt{cookies}.
\end{itemize}

\section{Fase 2: Diseño Arquitectónico y Prototipado}
En esta etapa se diseñará la solución técnica que operará dentro de las restricciones de seguridad del navegador.
\begin{itemize}
    \item \textbf{Diseño de Arquitectura de Software:} Se elaborarán diagramas de componentes y secuencia (UML) detallando la interacción entre el \textit{Service Worker} de la herramienta propuesta y la API \texttt{chrome.management}.
    \item \textbf{Diseño de Interfaz de Usuario (UI/UX):} Se diseñarán los prototipos de las alertas y el panel de control, priorizando la claridad para usuarios no técnicos.
    \item \textbf{Selección de herramientas:} Configuración del entorno de desarrollo utilizando tecnologías web estándar (TypeScript, React/Vue, Vite) y herramientas de empaquetado compatibles con Manifest V3.
\end{itemize}

\section{Fase 3: Construcción e Implementación del Mecanismo}
Fase central del proyecto donde se codificarán los módulos operativos de la extensión de seguridad.
\begin{itemize}
    \item \textbf{Módulo de Auditoría de Instalación:} Desarrollo del script que analiza el archivo \texttt{manifest.json} de otras extensiones al momento de su instalación y calcula un puntaje de riesgo inicial.
    \item \textbf{Motor de Detección de Cambios (Update Watcher):} Implementación de un sistema de almacenamiento local (\texttt{chrome.storage}) que guarda una "huella digital" de los permisos de cada extensión. El sistema comparará la versión $N$ vs $N+1$ tras cada actualización automática para alertar sobre elevación de privilegios.
    \item \textbf{Sistema de Control Contextual:} Desarrollo de \textit{listeners} que detectan la URL activa en la pestaña del usuario. Si la URL coincide con una lista de sitios críticos definidos por el usuario, el sistema desactivará temporalmente las extensiones no esenciales o bloqueará su inyección de scripts.
\end{itemize}

\section{Fase 4: Validación y Pruebas Controladas}
Esta fase busca verificar la eficacia técnica y la usabilidad de la herramienta.
\begin{itemize}
    \item \textbf{Creación de Escenarios de Ataque (Sandboxing):} Desarrollo de extensiones "simulacro" (dummies) que repliquen comportamientos maliciosos (ej. solicitar permiso de \texttt{cookies} tras una actualización) para probar la capacidad de detección de la herramienta.
    \item \textbf{Pruebas de Rendimiento:} Medición del impacto de la herramienta en la carga de páginas y uso de memoria RAM del navegador.
\end{itemize}

\endgroup

% ============================================
% CAPÍTULO 4: MODELO DE EVALUACIÓN DE RIESGO (FASE 1)
% ============================================

\chapter{Modelo de Evaluación de Riesgo}
\label{chap:modelo_riesgo}

\begingroup
\setlength{\parindent}{0pt}
\setlength{\parskip}{0.6em}

Este capítulo presenta el modelo de evaluación de riesgo desarrollado como resultado de la Fase 1 del proyecto. El modelo establece las reglas lógicas que el sistema utiliza para determinar si una extensión de Google Chrome representa un riesgo de seguridad o privacidad para el usuario. La taxonomía de permisos, las matrices de correlación y la fórmula de riesgo se basan en la documentación oficial de Chrome y en extensiones distribuidas a través de la Chrome Web Store. Los permisos exclusivos de ChromeOS y entornos empresariales se documentan pero se excluyen del cálculo de riesgo, dado que no son aplicables en instalaciones estándar de Chrome en Windows, macOS o Linux.

\section{Fundamentos y Metodología de Clasificación}

\subsection{Fuentes Primarias}

La clasificación de riesgo de permisos presentada en este capítulo se fundamenta en una síntesis de múltiples fuentes:

\begin{enumerate}
    \item \textbf{Google Permission Risk Whitepaper (2019):} Documento oficial de Google Chrome Enterprise \& Education que clasifica los permisos en cuatro niveles: \textit{Highest}, \textit{High}, \textit{Medium} y \textit{Low} \cite{google_permission_risk}. Aunque este whitepaper constituye la única guía oficial disponible, data de julio de 2019 (Chrome 75), previo a la implementación completa de Manifest V3.
    
    \item \textbf{LayerX Enterprise Browser Extension Security Report 2025:} Reporte de industria (no sometido a revisión por pares) publicado por LayerX Security, un proveedor de soluciones de seguridad de navegador empresarial. Combina estadísticas de uso real en entornos empresariales con análisis de permisos, basado en telemetría de clientes empresariales de LayerX. Encontró que el 53\% de usuarios empresariales tienen extensiones con permisos de riesgo ``High'' o ``Critical''. El reporte analiza la prevalencia de APIs sensibles como \texttt{cookies} (4.87\% de extensiones), \texttt{scripting} (14.34\%), \texttt{identity} (3.51\%) y \texttt{tabs} (23.92\%), y describe cualitativamente los riesgos asociados a cada una (robo de sesiones, inyección de código, exfiltración de credenciales). Nótese que LayerX clasifica extensiones por nivel de riesgo agregado, no permisos individuales. Los datos se utilizan en este trabajo como referencia complementaria a las fuentes académicas, reconociendo el posible sesgo inherente a un reporte con interés comercial \cite{layerx2025report}.

\needspace{8\baselineskip}

    
    \item \textbf{Investigación académica reciente:} El estudio de Singh et al. (2025) \cite{singh2025study} experimentalmente demostró que extensiones maliciosas pueden ser desarrolladas, publicadas y ejecutadas en la Chrome Web Store y Mozilla Add-ons Store usando permisos mínimos como \texttt{activeTab}, \texttt{scripting} y \texttt{storage}. Nayak et al. (2024) \cite{nayak2024experimental} documentan que el 12.5\% de las extensiones analizadas pueden extraer datos sensibles.
\end{enumerate}

\subsection{Criterios de Actualización}

Dado que el whitepaper de Google no ha sido actualizado desde 2019, este trabajo realiza ajustes basados en evidencia empírica de ataques documentados en 2024-2025. Los criterios para elevar el nivel de riesgo de un permiso respecto a la clasificación original de Google son:

\begin{itemize}
    \item \textbf{Uso documentado en campañas maliciosas:} Si un permiso ha sido explotado en ataques con impacto significativo (>100,000 usuarios afectados).
    \item \textbf{Cambio de funcionalidad en MV3:} Permisos cuyo comportamiento cambió significativamente con la transición a Manifest V3.
    \item \textbf{Consenso de múltiples fuentes:} Si el análisis de riesgos de LayerX y/o investigadores académicos (Singh et al., Nayak et al.) documentan el abuso de un permiso en contextos que justifican una clasificación más alta que la asignada por Google.
\end{itemize}

\section{Taxonomía de Riesgo de Permisos}

La clasificación que sigue cubre los \textbf{permisos declarables} en el campo \texttt{"permissions"} del \texttt{manifest.json}. Es importante distinguir cuatro mecanismos de acceso a funcionalidad en Chrome Extensions:

\begin{enumerate}
    \item \textbf{Permisos explícitos} (\texttt{"permissions"}): Strings declarados en el manifest que habilitan APIs específicas. Son el foco principal de esta taxonomía.
    \item \textbf{Permisos de host} (\texttt{"host\_permissions"}): Patrones de coincidencia (match patterns) que determinan sobre qué dominios puede operar la extensión. Incluyen patrones como \texttt{"https://example.com/*"}, \texttt{"*://*/*"} o \texttt{"<all\_urls>"}. En MV3 se declaran en un campo separado del manifest. Son analizados por el motor de auditoría como factor de alcance.
    \item \textbf{Claves del manifest}: Algunas APIs se activan mediante claves especiales del manifest, no mediante permisos. Ejemplos: \texttt{"action"} (icono de extensión), \texttt{"commands"} (atajos de teclado), \texttt{"omnibox"} (barra de búsqueda), \texttt{"devtools\_page"} (herramientas de desarrollo).
    \item \textbf{APIs sin permiso}: Varias APIs están disponibles sin declarar nada: \texttt{chrome.runtime} (excepto \texttt{connectNative}), funciones básicas de \texttt{chrome.tabs} y \texttt{chrome.windows} (sin campos sensibles), \texttt{chrome.i18n}, \texttt{chrome.extension} y \texttt{chrome.dom}.
\end{enumerate}

Para el modelo de riesgo propuesto, los permisos explícitos (categoría 1) y los permisos de host (categoría 2) son relevantes, ya que ambos aparecen en el manifest analizable de cada extensión y determinan las capacidades que la herramienta de auditoría puede evaluar.

\subsection{Nivel CRÍTICO}

Permisos que otorgan acceso irrestricto a datos del usuario o capacidad de escapar del sandbox del navegador. Google los clasifica como \textit{Highest Risk} \cite{google_permission_risk}.

\textbf{Nota sobre \texttt{host\_permissions}:} Los patrones de acceso a sitios web (\texttt{<all\_urls>}, \texttt{*://*/*}, dominios específicos) no son permisos del campo \texttt{"permissions"} del manifest, sino patrones declarados en el campo \texttt{"host\_permissions"}. Google los clasifica como el riesgo más alto porque determinan el alcance de acceso a sitios web \cite{google_permission_risk}. En el modelo de riesgo propuesto, estos patrones se incorporan exclusivamente mediante el \textbf{factor de alcance $f(H)$} (Sección \ref{sec:formula}), no como permisos individuales, para evitar un doble conteo entre la capacidad y el alcance de la extensión. El motor de auditoría analiza ambos campos del manifest.

\begin{table}[!htbp]
\centering
\small
\begin{tabularx}{\textwidth}{|p{3.5cm}|X|}
\hline
\textbf{Permiso} & \textbf{Riesgo y Justificación} \\
\hline
\texttt{debugger} & Acceso al backend del depurador de Chrome. Google lo describe como acceso de ``superusuario'' porque puede inspeccionar y modificar cualquier aspecto del navegador \cite{google_permission_risk}. \\
\hline
\texttt{nativeMessaging} & Comunicación directa con programas nativos del sistema operativo. Permite escapar completamente del sandbox del navegador \cite{google_permission_risk}. \textbf{Nota:} El permiso \texttt{runtime} de la lista oficial habilita las mismas funciones (\texttt{connectNative()} y \texttt{sendNativeMessage()}); ambos permisos otorgan capacidad equivalente. \\
\hline
\texttt{proxy} & Control completo de la configuración del proxy. Puede redirigir todo el tráfico del usuario a servidores controlados por el atacante \cite{google_permission_risk}. \\
\hline
\texttt{vpnProvider} & Crear túneles VPN y enviar/recibir paquetes a nivel de red. Acceso equivalente a \texttt{proxy} pero a nivel más bajo \cite{google_permission_risk}. \textbf{ChromeOS only.} \\
\hline
\texttt{tabCapture} & Captura de video/audio de la pestaña activa. Usado en ataques documentados para grabar videollamadas de Google Meet y Zoom sin permiso adicional \cite{singh2025study}. \\
\hline
\texttt{pageCapture} & Obtiene MHTML completo de cualquier página, incluyendo datos sensibles visibles como formularios y tokens \cite{google_permission_risk}. \\
\hline
\end{tabularx}
\caption{Permisos de nivel CRÍTICO}
\label{tab:critical_permissions}
\end{table}

\needspace{10\baselineskip}

\subsection{Nivel ALTO}

Permisos que permiten acceso significativo a datos del usuario o modificación del entorno de navegación. Google los clasifica como \textit{High Risk} \cite{google_permission_risk}. Incluye permisos que este trabajo eleva de nivel basándose en evidencia de ataques recientes.

\begingroup
\small
\begin{xltabular}{\textwidth}{|p{4.2cm}|X|}
\hline
\textbf{Permiso} & \textbf{Riesgo y Justificación} \\
\hline
\endfirsthead

\hline
\textbf{Permiso} & \textbf{Riesgo y Justificación} \\
\hline
\endhead

\hline
\multicolumn{2}{|r|}{\textit{Continúa en la siguiente página}} \\
\hline
\endfoot

\hline
\caption{Permisos de nivel ALTO}
\label{tab:high_permissions}
\endlastfoot

\texttt{cookies} & Google lo clasifica como \textit{Medium}. LayerX advierte que una extensión maliciosa con acceso a cookies podría ``steal or delete session cookies, effectively hijacking or terminating the user's session'' \cite{layerx2025report}. La campaña DataByCloud (enero 2026) exfiltró cookies de autenticación cada 60 segundos usando este permiso \cite{pandya2026databycloud}. \textbf{Elevado de Medium a Alto.} \textbf{Nota de dependencia funcional:} el permiso \texttt{cookies} es una \textit{capacidad} que, por sí sola, solo permite acceder a cookies del dominio de la propia extensión. Para acceder a cookies de dominios externos, la extensión debe poseer \texttt{host\_permissions} o \texttt{activeTab} para dichos dominios. Sin \texttt{host\_permissions}, el riesgo efectivo de \texttt{cookies} es mínimo. Se clasifica como Alto considerando el escenario combinado (cookies + host\_permissions), que es el patrón observado en ataques reales. \\
\hline
\texttt{scripting} & Google lo clasifica como \textit{Medium}. LayerX advierte que ``permite la inyección de código JavaScript en páginas web, que puede usarse para capturar formularios de login, pulsaciones de teclas, o manipular páginas web para exfiltrar credenciales'' \cite{layerx2025report}. Singh et al. demostraron que ``permisos mínimos---activeTab, scripting y storage---son suficientes para ejecutar acciones dañinas'' incluyendo keyloggers \cite{singh2025study}. \textbf{Elevado de Medium a Alto.} \textbf{Nota de dependencia funcional:} \texttt{scripting} es una \textit{capacidad} que habilita \texttt{chrome.scripting.executeScript()}, pero esta API requiere que la extensión tenga permisos de host (\texttt{host\_permissions} o \texttt{activeTab}) sobre la pestaña objetivo para ejecutarse. Sin acceso a hosts, \texttt{scripting} es inerte. La clasificación Alto refleja el escenario combinado con permisos de host, que es el patrón de abuso documentado. \\
\hline
\texttt{declarativeNetRequest} & Google lo clasifica como \textit{High}. Permite modificar peticiones de red mediante reglas declarativas, incluyendo la acción \texttt{modifyHeaders} que puede alterar cabeceras HTTP (cookies, referer, etc.) en peticiones salientes. Campañas como DataByCloud solicitaron este permiso en combinación con \texttt{cookies} y \texttt{host\_permissions} para maximizar su superficie de ataque \cite{pandya2026databycloud}. \textbf{Nota de dependencia funcional:} la acción \texttt{modifyHeaders} requiere \texttt{host\_permissions} o \texttt{declarativeNetRequestWithHostAccess} para los dominios objetivo. Las reglas de bloqueo básicas (\texttt{block}, \texttt{upgradeScheme}) no requieren permisos de host adicionales, pero las reglas de redirección (\texttt{redirect}) sí los requieren. \\
\hline
\texttt{webRequest} + \texttt{webRequestBlocking} & Ver cualquier petición GET/POST y sus URLs. Puede bloquear o modificar peticiones en tránsito. Google los clasifica como \textit{Medium} \cite{google_permission_risk}; este trabajo los eleva a Alto dado que permiten inspeccionar y modificar tráfico HTTP completo, incluyendo cabeceras de autenticación. \textbf{Nota:} En MV3, \texttt{webRequestBlocking} solo está disponible para extensiones instaladas por política empresarial; extensiones regulares deben usar \texttt{declarativeNetRequest}. \\
\hline
\texttt{history} & Acceso completo al historial de navegación: leer, añadir y eliminar entradas \cite{google_permission_risk}. \\
\hline
\texttt{downloads} + \texttt{downloads.open} & Iniciar, monitorear y manipular descargas. Con \texttt{downloads.open} puede ejecutar archivos descargados \cite{google_permission_risk}. \\
\hline
\texttt{privacy} & Puede desactivar protecciones como Safe Browsing, exponiendo al usuario a malware \cite{google_permission_risk}. \\
\hline
\texttt{browsingData} & Borrar datos de navegación. Útil para atacantes que quieren eliminar evidencia forense \cite{google_permission_risk}. \\
\hline
\texttt{contentSettings} & Puede permitir que plugins se ejecuten sin sandbox o modificar permisos de sitios \cite{google_permission_risk}. \\
\hline
\texttt{webNavigation} & Escuchar todos los sitios web que visita el usuario \cite{google_permission_risk}. \\
\hline
\texttt{userScripts} & Inyectar scripts de usuario en páginas web. Permiso nuevo en MV3 que requiere habilitación explícita del usuario \cite{chrome_permissions_list}. Similar en riesgo a \texttt{scripting}. \\
\hline
\texttt{webAuthenticationProxy} & Proxy de autenticación web. Permite interceptar flujos de autenticación WebAuthn \cite{chrome_permissions_list}. \\
\hline
\texttt{declarativeNetRequest\-WithHostAccess} & Variante de \texttt{declarativeNetRequest} que requiere \texttt{host\_permissions} explícitos para \textit{todas} las acciones, incluyendo bloqueo y redirección. No muestra warning de instalación propio, pero depende completamente de los hosts declarados \cite{chrome_permissions_list}. \textbf{Nota de dependencia funcional:} sin \texttt{host\_permissions}, este permiso es completamente inerte. \\
\hline
\texttt{certificateProvider} & Proveer certificados para autenticación TLS. Google lo clasifica como \textit{Low} \cite{google_permission_risk}, pero en contextos empresariales donde se gestiona autenticación mutua TLS, el compromiso de certificados puede habilitar ataques de suplantación. \textbf{ChromeOS only.} \\
\hline
\texttt{platformKeys} & Acceso a claves de plataforma para autenticación. Google lo clasifica como \textit{Low} \cite{google_permission_risk}, pero este trabajo lo eleva a Alto para entornos empresariales ChromeOS donde las claves de plataforma protegen autenticación corporativa. \textbf{ChromeOS only.} \\
\hline
\texttt{desktopCapture} & Captura de pantalla completa del escritorio, incluyendo aplicaciones fuera del navegador. Google lo clasifica como \textit{Medium} \cite{google_permission_risk}. Este trabajo lo eleva a Alto dado que permite capturar información de cualquier aplicación visible en pantalla, incluyendo ventanas de correo, gestores de contraseñas y terminales, lo que excede el alcance del propio navegador. \textbf{Elevado de Medium a Alto.} \\
\hline
\end{xltabular}
\endgroup

\needspace{10\baselineskip}
\subsection{Nivel MEDIO}

Permisos que representan riesgo moderado y requieren evaluación contextual. Google los clasifica como \textit{Medium Risk} \cite{google_permission_risk}. Este trabajo eleva \texttt{alarms} de \textit{Low} a \textit{Medium} basándose en su uso como vector de persistencia en MV3.

{\small
\begin{longtable}{|p{3.5cm}|p{11cm}|}
\hline
\textbf{Permiso} & \textbf{Riesgo y Justificación} \\
\hline
\endfirsthead
\hline
\endhead
\endfoot
\hline
\caption{Permisos de nivel MEDIO}
\label{tab:medium_permissions}
\endlastfoot
\texttt{alarms} & Google lo clasifica como \textit{Low} (``establecer temporizadores''). Sin embargo, en MV3 este permiso se ha convertido en el mecanismo dominante para mantener Service Workers activos. Aunque existen otros mecanismos de persistencia (como \texttt{connectNative}), estos requieren componentes externos que elevan la barrera para atacantes. En contraste, \texttt{chrome.alarms} permite crear ``heartbeats'' que despiertan al worker periódicamente sin dependencias adicionales \cite{chrome_alarms_docs}. La campaña Phantom Shuttle transmitía credenciales robadas ``cada cinco minutos'' mediante este mecanismo \cite{socket2025phantom}. \textbf{Elevado de Low a Medio.} \\
\hline
\texttt{activeTab} & Acceso temporal a la pestaña activa solo tras interacción del usuario (clic en icono). Google lo considera más seguro que host\_permissions permanentes \cite{google_permission_risk}. \\
\hline
\texttt{bookmarks} & Crear, organizar y manipular marcadores del usuario \cite{google_permission_risk}. \\
\hline
\texttt{clipboardRead} / \texttt{clipboardWrite} & Leer o escribir en el portapapeles del sistema \cite{google_permission_risk}. \\
\hline
\texttt{geolocation} & Permite usar la API de geolocalización sin que el navegador muestre el prompt estándar de permiso al usuario. En MV3, \texttt{navigator.geolocation} no está disponible directamente en el Service Worker (no tiene DOM); debe invocarse desde un content script o un documento \texttt{offscreen} \cite{google_permission_risk}. \\
\hline
\texttt{identity} / \texttt{identity.email} & Solicitar OAuth con credenciales del usuario u obtener su email \cite{google_permission_risk}. \\
\hline
\texttt{management} & Gestionar la lista de extensiones/apps instaladas. Podría usarse para deshabilitar extensiones de seguridad \cite{google_permission_risk}. \\
\hline
\texttt{sessions} & Consultar y restaurar pestañas/ventanas de sesiones anteriores \cite{google_permission_risk}. \\
\hline
\texttt{topSites} & Obtener o modificar la lista de sitios más visitados \cite{google_permission_risk}. \\
\hline
\texttt{contextMenus} & Integrar menús de la extensión en el menú contextual del navegador \cite{google_permission_risk}. \\
\hline
\texttt{tabGroups} & Ver y gestionar grupos de pestañas del usuario. Permiso nuevo en MV3 \cite{chrome_permissions_list}. \\
\hline
\texttt{dns} & Acceso a la API de resolución DNS del navegador \cite{chrome_permissions_list}. \\
\hline
\texttt{tabs} & No otorga acceso al namespace \texttt{chrome.tabs} (que no requiere permiso), sino a campos sensibles de los objetos Tab: \texttt{url}, \texttt{pendingUrl}, \texttt{title} y \texttt{favIconUrl}. Warning: ``Read your browsing history'' \cite{chrome_permissions_list}. Google lo clasifica como \textit{High} \cite{google_permission_risk}; este trabajo lo ubica en Medio con la siguiente justificación: (1) a diferencia de \texttt{history}, solo expone pestañas activas (no el historial completo) y no permite modificar contenido; (2) la revisión de campañas maliciosas documentadas en 2024--2026 (DarkSpectre, DataByCloud, Phantom Shuttle) no identifica a \texttt{tabs} como permiso clave explotado, a diferencia de \texttt{cookies} y \texttt{scripting}; (3) con el 23.92\% de extensiones solicitando este permiso \cite{layerx2025report}, mantenerlo como Alto generaría un volumen excesivo de alertas que diluiría la señal del modelo. No obstante, sigue siendo un vector de tracking significativo. \textbf{Rebajado de High a Medio.} \\
\hline
\texttt{offscreen} & Crear documentos offscreen con acceso completo al DOM, compensando las limitaciones de los Service Workers en MV3. Un documento offscreen puede ejecutar \texttt{setInterval}, mantener WebSockets abiertos y acceder a APIs que requieren DOM (como \texttt{navigator.geolocation}) \cite{chrome_permissions_list}. Representa un vector de persistencia y exfiltración en MV3. \textbf{Clasificado como Medio (sin clasificación previa de Google, ya que este permiso fue introducido en MV3, posterior al whitepaper de 2019).} \\
\hline
\texttt{processes} & Monitorear procesos del navegador, obtener métricas de CPU/memoria, y terminar procesos. Disponible solo en canal Dev \cite{chrome_permissions_list}. \\
\hline
\end{longtable}
}

\subsection{Nivel BAJO}

Permisos que representan riesgo mínimo para la seguridad del usuario. Google los clasifica como \textit{Low Risk} \cite{google_permission_risk}. Se incluyen permisos nuevos introducidos en Manifest V3 que no tienen clasificación oficial.

{\small
\begin{longtable}{|p{4.5cm}|p{11cm}|}
\hline
\textbf{Permiso} & \textbf{Riesgo y Justificación} \\
\hline
\endfirsthead
\hline
\textbf{Permiso} & \textbf{Riesgo y Justificación} \\
\hline
\endhead
\hline
\multicolumn{2}{|r|}{\textit{Continúa en la siguiente página}} \\
\hline
\endfoot
\hline
\caption{Permisos de nivel BAJO}
\label{tab:low_permissions}
\endlastfoot
\texttt{storage} / \texttt{unlimitedStorage} & Almacenar datos en el computador del usuario. No permite acceso a datos de otras extensiones o sitios web \cite{google_permission_risk}. \\
\hline
\texttt{notifications} & Mostrar notificaciones popup al usuario en el escritorio \cite{google_permission_risk}. \\
\hline
\texttt{idle} & Determinar si el sistema está siendo usado activamente \cite{google_permission_risk}. \\
\hline
\texttt{power} & Gestionar estado de energía (prevenir suspensión) \cite{google_permission_risk}. \\
\hline
\texttt{tts} / \texttt{ttsEngine} & API de texto a voz para accesibilidad \cite{google_permission_risk}. \\
\hline
\texttt{fontSettings} & Gestionar configuración de fuentes del navegador \cite{google_permission_risk}. \\
\hline
\texttt{declarativeContent} & Habilitar acciones según contenido de página sin leer el contenido \cite{google_permission_risk}. \\
\hline
\texttt{gcm} & Enviar/recibir mensajes vía Google Cloud Messaging \cite{google_permission_risk}. \\
\hline
\texttt{sidePanel} & Mostrar contenido en el panel lateral de Chrome. Permiso nuevo en MV3 \cite{chrome_permissions_list}. \\
\hline
\texttt{search} & Interactuar con la barra de búsqueda de Chrome \cite{chrome_permissions_list}. \\
\hline
\texttt{favicon} & Leer íconos de sitios web visitados. Permiso nuevo en MV3 \cite{chrome_permissions_list}. \\
\hline
\texttt{readingList} & Leer y modificar la lista de lectura del usuario \cite{chrome_permissions_list}. \\
\hline
\texttt{printing} / \texttt{printingMetrics} & Acceso a funciones de impresión del navegador \cite{chrome_permissions_list}. \textbf{ChromeOS only.} \\
\hline
\texttt{documentScan} & Acceso a escáneres de documentos \cite{chrome_permissions_list}. \textbf{ChromeOS only.} \\
\hline
\texttt{loginState} & Detectar estado de login \cite{chrome_permissions_list}. \textbf{ChromeOS only.} \\
\hline
\texttt{accessibilityFeatures.*} & Leer y modificar características de accesibilidad del navegador (ej. alto contraste, lupa). Riesgo mínimo ya que solo afecta configuraciones de UI del navegador \cite{chrome_permissions_list}. \\
\hline
\texttt{background} & Permite que Chrome arranque temprano (al iniciar sesión del SO) y no cierre al cerrar la última ventana. No otorga acceso a datos del usuario \cite{chrome_permissions_list}. \\
\hline
\texttt{declarativeNetRequest\-Feedback} & Permite escribir errores y warnings en la consola de DevTools para depuración de reglas de \texttt{declarativeNetRequest}. Solo funciona con extensiones desempaquetadas; ignorado para extensiones del Chrome Web Store \cite{chrome_permissions_list}. \\
\hline
\texttt{downloads.ui} & Permite modificar la interfaz de usuario de descargas de Chrome mediante \texttt{chrome.downloads.setUiOptions()} \cite{chrome_permissions_list}. \\
\hline
\texttt{system.cpu} & Obtener información sobre el CPU del sistema \cite{chrome_permissions_list}. \\
\hline
\texttt{system.display} & Obtener información sobre las pantallas conectadas al sistema \cite{chrome_permissions_list}. \\
\hline
\texttt{system.memory} & Obtener información sobre la memoria RAM del sistema \cite{chrome_permissions_list}. \\
\hline
\texttt{system.storage} & Obtener información sobre dispositivos de almacenamiento. Puede expulsar dispositivos removibles. Google lo clasifica como \textit{Medium} \cite{google_permission_risk}; este trabajo lo ubica en Bajo porque: (1) la capacidad de expulsar dispositivos tiene impacto limitado en el contexto de privacidad de extensiones web; (2) la información de almacenamiento expuesta (nombre, tipo, capacidad) no incluye datos personales del usuario; (3) no se documentan ataques que exploten este permiso en campañas de extensiones maliciosas. \textbf{Rebajado de Medium a Bajo.} \\
\hline
\texttt{printerProvider} & Implementar un proveedor de impresora para que Chrome reconozca la extensión como destino de impresión \cite{chrome_permissions_list}. \\
\hline
\end{longtable}
}

\subsection{Permisos Exclusivos de ChromeOS y Entornos Empresariales}
\label{sec:chromeos_permissions}

La documentación oficial de Chrome incluye permisos que son \textbf{exclusivos de ChromeOS} o que requieren \textbf{políticas empresariales} (Enterprise) para funcionar. Estos permisos no están disponibles para extensiones instaladas en Windows, macOS o Linux estándar, por lo que el motor de auditoría propuesto debe identificarlos y excluirlos del cálculo de riesgo cuando la extensión objetivo no opera en dichos entornos.

Los permisos ya clasificados en las tablas anteriores que son exclusivos de ChromeOS se marcaron con la nota \textbf{ChromeOS only}.

Adicionalmente, los siguientes permisos de ChromeOS/Enterprise no fueron incluidos en la clasificación principal por no ser relevantes para el modelo de riesgo en navegadores de escritorio estándar:

\begin{table}[H]
\centering
\small
\begin{tabularx}{\textwidth}{|p{5.5cm}|X|}
\hline
\textbf{Permiso} & \textbf{Descripción y Plataforma} \\
\hline
\texttt{audio} & Control de dispositivos de audio del sistema (volumen, mute, dispositivos activos). \textbf{ChromeOS only} (modo kiosco) \cite{chrome_permissions_list}. \\
\hline
\texttt{enterprise.deviceAttributes} & Obtener atributos del dispositivo empresarial. \textbf{Enterprise policy required} \cite{chrome_permissions_list}. \\
\hline
\texttt{enterprise.hardwarePlatform} & Información del hardware de la plataforma empresarial. \textbf{Enterprise policy required} \cite{chrome_permissions_list}. \\
\hline
\texttt{enterprise.networkingAttributes} & Atributos de red del dispositivo empresarial. \textbf{Enterprise policy required} \cite{chrome_permissions_list}. \\
\hline
\texttt{enterprise.platformKeys} & Claves de plataforma para autenticación empresarial. \textbf{Enterprise policy required} \cite{chrome_permissions_list}. \\
\hline
\texttt{fileBrowserHandler} & Manejador del explorador de archivos de ChromeOS. \textbf{ChromeOS only} \cite{chrome_permissions_list}. \\
\hline
\texttt{fileSystemProvider} & Proveer sistemas de archivos virtuales al explorador de ChromeOS. \textbf{ChromeOS only} \cite{chrome_permissions_list}. \\
\hline
\texttt{input} & Editor de métodos de entrada (IME) personalizado. \textbf{ChromeOS only} \cite{chrome_permissions_list}. \\
\hline
\texttt{wallpaper} & Cambiar el fondo de pantalla del dispositivo. \textbf{ChromeOS only} \cite{chrome_permissions_list}. \\
\hline
\end{tabularx}
\caption{Permisos exclusivos de ChromeOS y entornos empresariales (excluidos del modelo de riesgo)}
\label{tab:chromeos_permissions}
\end{table}

\subsection{Resumen de Diferencias con el Whitepaper de Google}

\begin{table}[H]
\centering
\small
\begin{tabular}{|l|c|c|p{6cm}|}
\hline
\textbf{Permiso} & \textbf{Google (2019)} & \textbf{Este trabajo} & \textbf{Justificación del cambio} \\
\hline
\texttt{cookies} & Medium & \textbf{Alto} & Explotado en campaña DataByCloud para session hijacking (2026) \cite{pandya2026databycloud}. LayerX documenta que extensiones con acceso a cookies pueden secuestrar sesiones \cite{layerx2025report}. \\
\hline
\texttt{scripting} & Medium & \textbf{Alto} & Usado para keyloggers en MV3 demostrado experimentalmente por Singh et al. (2025) \cite{singh2025study}. LayerX documenta riesgo de captura de formularios y exfiltración de credenciales \cite{layerx2025report}. \\
\hline
\texttt{alarms} & Low & \textbf{Medio} & Vector de persistencia estándar en MV3 para mantener Service Workers activos. Usado en heartbeat de exfiltración \cite{socket2025phantom}. \\
\hline
\texttt{offscreen} & N/A (MV3) & \textbf{Medio} & Permiso nuevo en MV3, sin clasificación previa de Google. Permite crear documentos ocultos con DOM completo, ejecutar timers persistentes y WebSockets. Vector de persistencia y exfiltración en MV3. \\
\hline
\texttt{desktopCapture} & Medium & \textbf{Alto} & Permite capturar información de cualquier aplicación visible en pantalla, incluyendo ventanas fuera del navegador (correo, gestores de contraseñas, terminales), excediendo el alcance del propio navegador. \\
\hline
\texttt{tabs} & High & \textbf{Medio} & Solo expone pestañas activas (no historial completo), no permite modificar contenido, y no se documenta como vector clave en campañas maliciosas recientes. Su alta prevalencia (23.92\%) generaría exceso de alertas. \\
\hline
\texttt{system.storage} & Medium & \textbf{Bajo} & Impacto limitado en privacidad: la información expuesta no incluye datos personales y no se documentan ataques explotando este permiso. \\
\hline
\end{tabular}
\caption{Permisos con clasificación actualizada respecto al whitepaper de Google}
\label{tab:permission_diff}
\end{table}


\section{Matrices de Correlación: Categorías vs. Permisos}

Las siguientes matrices constituyen una \textbf{contribución original} de este proyecto. Establecen la correlación entre las 18 categorías de extensiones disponibles en la Chrome Web Store\footnote{Categorías obtenidas de \url{https://chromewebstore.google.com/category/extensions}, consultado en marzo de 2026. Las categorías han sido reorganizadas por Google en tres grupos: Productivity, Lifestyle y Make Chrome Yours.} y los permisos esperados, sospechosos e incoherentes para cada tipo de extensión. Google no proporciona esta información en su documentación oficial.

Los permisos se clasifican según su coherencia con la funcionalidad declarada:
\begin{itemize}
    \item \textbf{Esperados} (verde): Permisos coherentes con la funcionalidad típica de la categoría.
    \item \textbf{Sospechosos} (amarillo): Permisos que requieren justificación adicional pero pueden ser legítimos.
    \item \textbf{Incoherentes} (rojo): Permisos sin justificación razonable que deben generar alertas.
\end{itemize}

\vspace{0.5em}

\subsection{Categorías de Productividad}

\noindent
\begin{table}[!htbp]
\centering
\small
\begin{tabular}{|p{2.5cm}|p{11.5cm}|}
\hline
\rowcolor{headerBlue}
\multicolumn{2}{|c|}{\textcolor{white}{\textbf{Developer Tools (Herramientas de Desarrollo)}}} \\
\hline
\rowcolor{expectedGreen}
\textbf{Esperados} & \texttt{storage}, \texttt{tabs}, \texttt{activeTab}, \texttt{scripting} (limitado), \texttt{contextMenus}, \texttt{notifications}, \texttt{sidePanel} \\
\hline
\rowcolor{suspiciousYellow}
\textbf{Sospechosos} & \texttt{cookies}, \texttt{history}, \texttt{webNavigation}, \texttt{downloads}, \texttt{clipboardRead} \\
\hline
\rowcolor{incoherentRed}
\textbf{Incoherentes} & \texttt{debugger} (excepto depuradores), \texttt{nativeMessaging}, \texttt{proxy}, \texttt{desktopCapture}, \texttt{tabCapture}, \texttt{<all\_urls>} \\
\hline
\end{tabular}
\end{table}

\noindent
\begin{table}[!htbp]
\centering
\small
\begin{tabular}{|p{2.5cm}|p{11.5cm}|}
\hline
\rowcolor{headerBlue}
\multicolumn{2}{|c|}{\textcolor{white}{\textbf{Workflow \& Planning (Flujo de Trabajo)}}} \\
\hline
\rowcolor{expectedGreen}
\textbf{Esperados} & \texttt{storage}, \texttt{activeTab}, \texttt{notifications}, \texttt{alarms}, \texttt{identity}, \texttt{contextMenus}, \texttt{sidePanel} \\
\hline
\rowcolor{suspiciousYellow}
\textbf{Sospechosos} & \texttt{tabs} (con URLs), \texttt{history}, \texttt{bookmarks}, \texttt{clipboardRead}, \texttt{scripting} (amplio) \\
\hline
\rowcolor{incoherentRed}
\textbf{Incoherentes} & \texttt{cookies}, \texttt{debugger}, \texttt{nativeMessaging}, \texttt{webRequest}, \texttt{proxy}, \texttt{desktopCapture}, \texttt{tabCapture} \\
\hline
\end{tabular}
\end{table}

\noindent
\begin{table}[!htbp]
\centering
\small
\begin{tabular}{|p{2.5cm}|p{11.5cm}|}
\hline
\rowcolor{headerBlue}
\multicolumn{2}{|c|}{\textcolor{white}{\textbf{Education (Educación)}}} \\
\hline
\rowcolor{expectedGreen}
\textbf{Esperados} & \texttt{storage}, \texttt{activeTab}, \texttt{tts}, \texttt{notifications}, \texttt{contextMenus}, \texttt{sidePanel} \\
\hline
\rowcolor{suspiciousYellow}
\textbf{Sospechosos} & \texttt{tabs}, \texttt{history}, \texttt{scripting} (amplio), \texttt{downloads}, \texttt{identity.email} \\
\hline
\rowcolor{incoherentRed}
\textbf{Incoherentes} & \texttt{cookies}, \texttt{debugger}, \texttt{nativeMessaging}, \texttt{webRequest}, \texttt{proxy}, \texttt{clipboardRead}, \texttt{desktopCapture} \\
\hline
\end{tabular}
\end{table}

\noindent
\begin{table}[!htbp]
\centering
\small
\begin{tabular}{|p{2.5cm}|p{11.5cm}|}
\hline
\rowcolor{headerBlue}
\multicolumn{2}{|c|}{\textcolor{white}{\textbf{Communication (Comunicación)}}} \\
\hline
\rowcolor{expectedGreen}
\textbf{Esperados} & \texttt{storage}, \texttt{activeTab}, \texttt{notifications}, \texttt{identity}, \texttt{contextMenus}, \texttt{sidePanel} \\
\hline
\rowcolor{suspiciousYellow}
\textbf{Sospechosos} & \texttt{tabs}, \texttt{scripting} (sitios específicos), \texttt{history}, \texttt{clipboardRead}, \texttt{clipboardWrite} \\
\hline
\rowcolor{incoherentRed}
\textbf{Incoherentes} & \texttt{cookies} (excepto propios), \texttt{debugger}, \texttt{nativeMessaging}, \texttt{proxy}, \texttt{webRequest}, \texttt{desktopCapture} \\
\hline
\end{tabular}
\end{table}

\begin{table}[H]
\centering
\small
\begin{tabular}{|p{2.5cm}|p{11.5cm}|}
\hline
\rowcolor{headerBlue}
\multicolumn{2}{|c|}{\textcolor{white}{\textbf{Tools (Herramientas Generales)}}} \\
\hline
\rowcolor{expectedGreen}
\textbf{Esperados} & \texttt{storage}, \texttt{activeTab}, \texttt{contextMenus}, \texttt{notifications}, \texttt{sidePanel}, \texttt{alarms} \\
\hline
\rowcolor{suspiciousYellow}
\textbf{Sospechosos} & \texttt{tabs}, \texttt{scripting}, \texttt{clipboardRead}, \texttt{clipboardWrite}, \texttt{downloads}, \texttt{identity}, \texttt{bookmarks}, \texttt{history} \\
\hline
\rowcolor{incoherentRed}
\textbf{Incoherentes} & \texttt{debugger}, \texttt{nativeMessaging}, \texttt{proxy}, \texttt{desktopCapture}, \texttt{tabCapture}, \texttt{webRequest}, \texttt{<all\_urls>} \\
\hline
\end{tabular}
\end{table}


\subsection{Categorías de Estilo de Vida}

\noindent
\begin{table}[!htbp]
\centering
\small
\begin{tabular}{|p{2.5cm}|p{11.5cm}|}
\hline
\rowcolor{headerBlue}
\multicolumn{2}{|c|}{\textcolor{white}{\textbf{Entertainment (Entretenimiento)}}} \\
\hline
\rowcolor{expectedGreen}
\textbf{Esperados} & \texttt{storage}, \texttt{activeTab}, \texttt{notifications}, \texttt{contextMenus} \\
\hline
\rowcolor{suspiciousYellow}
\textbf{Sospechosos} & \texttt{tabs}, \texttt{scripting}, \texttt{history}, \texttt{downloads}, \texttt{bookmarks} \\
\hline
\rowcolor{incoherentRed}
\textbf{Incoherentes} & \texttt{cookies}, \texttt{debugger}, \texttt{nativeMessaging}, \texttt{proxy}, \texttt{webRequest}, \texttt{identity}, \texttt{clipboardRead}, \texttt{geolocation} \\
\hline
\end{tabular}
\end{table}

\noindent
\begin{table}[!htbp]
\centering
\small
\begin{tabular}{|p{2.5cm}|p{11.5cm}|}
\hline
\rowcolor{headerBlue}
\multicolumn{2}{|c|}{\textcolor{white}{\textbf{Games (Juegos)}}} \\
\hline
\rowcolor{expectedGreen}
\textbf{Esperados} & \texttt{storage}, \texttt{notifications}, \texttt{alarms} \\
\hline
\rowcolor{suspiciousYellow}
\textbf{Sospechosos} & \texttt{activeTab}, \texttt{tabs}, \texttt{scripting}, \texttt{identity}, \texttt{geolocation} \\
\hline
\rowcolor{incoherentRed}
\textbf{Incoherentes} & \texttt{cookies}, \texttt{debugger}, \texttt{nativeMessaging}, \texttt{proxy}, \texttt{webRequest}, \texttt{history}, \texttt{downloads}, \texttt{clipboardRead} \\
\hline
\end{tabular}
\end{table}

\noindent
\begin{table}[!htbp]
\centering
\small
\begin{tabular}{|p{2.5cm}|p{11.5cm}|}
\hline
\rowcolor{headerBlue}
\multicolumn{2}{|c|}{\textcolor{white}{\textbf{News \& Weather (Noticias y Clima)}}} \\
\hline
\rowcolor{expectedGreen}
\textbf{Esperados} & \texttt{storage}, \texttt{notifications}, \texttt{alarms}, \texttt{geolocation} (para clima) \\
\hline
\rowcolor{suspiciousYellow}
\textbf{Sospechosos} & \texttt{activeTab}, \texttt{tabs}, \texttt{history}, \texttt{bookmarks}, \texttt{identity} \\
\hline
\rowcolor{incoherentRed}
\textbf{Incoherentes} & \texttt{cookies}, \texttt{debugger}, \texttt{nativeMessaging}, \texttt{proxy}, \texttt{webRequest}, \texttt{scripting} (amplio), \texttt{downloads} \\
\hline
\end{tabular}
\end{table}

\noindent
\begin{table}[!htbp]
\centering
\small
\begin{tabular}{|p{2.5cm}|p{11.5cm}|}
\hline
\rowcolor{headerBlue}
\multicolumn{2}{|c|}{\textcolor{white}{\textbf{Shopping (Compras)}}} \\
\hline
\rowcolor{expectedGreen}
\textbf{Esperados} & \texttt{storage}, \texttt{activeTab}, \texttt{notifications}, \texttt{contextMenus} \\
\hline
\rowcolor{suspiciousYellow}
\textbf{Sospechosos} & \texttt{tabs}, \texttt{scripting}, \texttt{cookies} (cupones), \texttt{history}, \texttt{clipboardRead}, \texttt{webNavigation} \\
\hline
\rowcolor{incoherentRed}
\textbf{Incoherentes} & \texttt{debugger}, \texttt{nativeMessaging}, \texttt{proxy}, \texttt{desktopCapture}, \texttt{tabCapture}, \texttt{downloads}, \texttt{identity} \\
\hline
\end{tabular}
\end{table}

\noindent
\begin{table}[!htbp]
\centering
\small
\begin{tabular}{|p{2.5cm}|p{11.5cm}|}
\hline
\rowcolor{headerBlue}
\multicolumn{2}{|c|}{\textcolor{white}{\textbf{Social Networking (Redes Sociales)}}} \\
\hline
\rowcolor{expectedGreen}
\textbf{Esperados} & \texttt{storage}, \texttt{activeTab}, \texttt{notifications}, \texttt{contextMenus}, \texttt{identity} \\
\hline
\rowcolor{suspiciousYellow}
\textbf{Sospechosos} & \texttt{tabs}, \texttt{scripting}, \texttt{history}, \texttt{downloads}, \texttt{clipboardRead}, \texttt{clipboardWrite} \\
\hline
\rowcolor{incoherentRed}
\textbf{Incoherentes} & \texttt{cookies} (excepto propios), \texttt{debugger}, \texttt{nativeMessaging}, \texttt{proxy}, \texttt{webRequest}, \texttt{<all\_urls>} \\
\hline
\end{tabular}
\end{table}

\noindent
\begin{table}[!htbp]
\centering
\small
\begin{tabular}{|p{2.5cm}|p{11.5cm}|}
\hline
\rowcolor{headerBlue}
\multicolumn{2}{|c|}{\textcolor{white}{\textbf{Travel (Viajes)}}} \\
\hline
\rowcolor{expectedGreen}
\textbf{Esperados} & \texttt{storage}, \texttt{notifications}, \texttt{geolocation}, \texttt{contextMenus} \\
\hline
\rowcolor{suspiciousYellow}
\textbf{Sospechosos} & \texttt{activeTab}, \texttt{tabs}, \texttt{scripting}, \texttt{history}, \texttt{bookmarks}, \texttt{identity} \\
\hline
\rowcolor{incoherentRed}
\textbf{Incoherentes} & \texttt{cookies}, \texttt{debugger}, \texttt{nativeMessaging}, \texttt{proxy}, \texttt{webRequest}, \texttt{downloads}, \texttt{clipboardRead} \\
\hline
\end{tabular}
\end{table}

\noindent
\begin{table}[!htbp]
\centering
\small
\begin{tabular}{|p{2.5cm}|p{11.5cm}|}
\hline
\rowcolor{headerBlue}
\multicolumn{2}{|c|}{\textcolor{white}{\textbf{Well-being (Bienestar)}}} \\
\hline
\rowcolor{expectedGreen}
\textbf{Esperados} & \texttt{storage}, \texttt{notifications}, \texttt{alarms}, \texttt{idle} \\
\hline
\rowcolor{suspiciousYellow}
\textbf{Sospechosos} & \texttt{activeTab}, \texttt{tabs}, \texttt{history}, \texttt{identity}, \texttt{geolocation} \\
\hline
\rowcolor{incoherentRed}
\textbf{Incoherentes} & \texttt{cookies}, \texttt{debugger}, \texttt{nativeMessaging}, \texttt{proxy}, \texttt{webRequest}, \texttt{scripting}, \texttt{downloads} \\
\hline
\end{tabular}
\end{table}

\begin{table}[H]
\centering
\small
\begin{tabular}{|p{2.5cm}|p{11.5cm}|}
\hline
\rowcolor{headerBlue}
\multicolumn{2}{|c|}{\textcolor{white}{\textbf{Household (Hogar)}}} \\
\hline
\rowcolor{expectedGreen}
\textbf{Esperados} & \texttt{storage}, \texttt{notifications}, \texttt{activeTab}, \texttt{contextMenus} \\
\hline
\rowcolor{suspiciousYellow}
\textbf{Sospechosos} & \texttt{tabs}, \texttt{scripting}, \texttt{clipboardRead}, \texttt{clipboardWrite}, \texttt{identity}, \texttt{geolocation} \\
\hline
\rowcolor{incoherentRed}
\textbf{Incoherentes} & \texttt{cookies}, \texttt{debugger}, \texttt{nativeMessaging}, \texttt{proxy}, \texttt{webRequest}, \texttt{desktopCapture}, \texttt{tabCapture}, \texttt{history} \\
\hline
\end{tabular}
\end{table}

\noindent
\begin{table}[!htbp]
\centering
\small
\begin{tabular}{|p{2.5cm}|p{11.5cm}|}
\hline
\rowcolor{headerBlue}
\multicolumn{2}{|c|}{\textcolor{white}{\textbf{Art \& Design (Arte y Diseño)}}} \\
\hline
\rowcolor{expectedGreen}
\textbf{Esperados} & \texttt{storage}, \texttt{activeTab}, \texttt{contextMenus}, \texttt{notifications} \\
\hline
\rowcolor{suspiciousYellow}
\textbf{Sospechosos} & \texttt{tabs}, \texttt{scripting}, \texttt{downloads}, \texttt{clipboardRead}, \texttt{clipboardWrite}, \texttt{desktopCapture} (captura diseño) \\
\hline
\rowcolor{incoherentRed}
\textbf{Incoherentes} & \texttt{cookies}, \texttt{debugger}, \texttt{nativeMessaging}, \texttt{proxy}, \texttt{webRequest}, \texttt{history}, \texttt{bookmarks}, \texttt{identity} \\
\hline
\end{tabular}
\end{table}

\noindent
\begin{table}[!htbp]
\centering
\small
\begin{tabular}{|p{2.5cm}|p{11.5cm}|}
\hline
\rowcolor{headerBlue}
\multicolumn{2}{|c|}{\textcolor{white}{\textbf{Just for Fun (Solo por Diversión)}}} \\
\hline
\rowcolor{expectedGreen}
\textbf{Esperados} & \texttt{storage}, \texttt{notifications}, \texttt{contextMenus} \\
\hline
\rowcolor{suspiciousYellow}
\textbf{Sospechosos} & \texttt{activeTab}, \texttt{tabs}, \texttt{scripting}, \texttt{alarms}, \texttt{identity} \\
\hline
\rowcolor{incoherentRed}
\textbf{Incoherentes} & \texttt{cookies}, \texttt{debugger}, \texttt{nativeMessaging}, \texttt{proxy}, \texttt{webRequest}, \texttt{history}, \texttt{downloads}, \texttt{clipboardRead}, \texttt{geolocation} \\
\hline
\end{tabular}
\end{table}

\subsection{Categorías de Personalizar Chrome}

\noindent
\begin{table}[!htbp]
\centering
\small
\begin{tabular}{|p{2.5cm}|p{11.5cm}|}
\hline
\rowcolor{headerBlue}
\multicolumn{2}{|c|}{\textcolor{white}{\textbf{Accessibility (Accesibilidad)}}} \\
\hline
\rowcolor{expectedGreen}
\textbf{Esperados} & \texttt{storage}, \texttt{activeTab}, \texttt{tts}, \texttt{accessibilityFeatures.*}, \texttt{fontSettings} \\
\hline
\rowcolor{suspiciousYellow}
\textbf{Sospechosos} & \texttt{tabs}, \texttt{scripting} (amplio), \texttt{history}, \texttt{contextMenus} \\
\hline
\rowcolor{incoherentRed}
\textbf{Incoherentes} & \texttt{cookies}, \texttt{debugger}, \texttt{nativeMessaging}, \texttt{proxy}, \texttt{webRequest}, \texttt{downloads}, \texttt{clipboardRead}, \texttt{identity} \\
\hline
\end{tabular}
\end{table}

\noindent
\begin{table}[!htbp]
\centering
\small
\begin{tabular}{|p{2.5cm}|p{11.5cm}|}
\hline
\rowcolor{headerBlue}
\multicolumn{2}{|c|}{\textcolor{white}{\textbf{Functionality \& UI (Funcionalidad e Interfaz)}}} \\
\hline
\rowcolor{expectedGreen}
\textbf{Esperados} & \texttt{storage}, \texttt{activeTab}, \texttt{contextMenus}, \texttt{sidePanel}, \texttt{notifications}, \texttt{tabs} (sin URLs) \\
\hline
\rowcolor{suspiciousYellow}
\textbf{Sospechosos} & \texttt{scripting}, \texttt{history}, \texttt{bookmarks}, \texttt{downloads}, \texttt{tabGroups} \\
\hline
\rowcolor{incoherentRed}
\textbf{Incoherentes} & \texttt{cookies}, \texttt{debugger}, \texttt{nativeMessaging}, \texttt{proxy}, \texttt{webRequest}, \texttt{clipboardRead}, \texttt{identity} \\
\hline
\end{tabular}
\end{table}

\noindent
\begin{table}[!htbp]
\centering
\small
\begin{tabular}{|p{2.5cm}|p{11.5cm}|}
\hline
\rowcolor{headerBlue}
\multicolumn{2}{|c|}{\textcolor{white}{\textbf{Privacy \& Security (Privacidad y Seguridad)}}} \\
\hline
\rowcolor{expectedGreen}
\textbf{Esperados} & \texttt{storage}, \texttt{activeTab}, \texttt{declarativeNetRequest}, \texttt{privacy}, \texttt{cookies} (limpieza), \texttt{browsingData}, \texttt{webRequest} (solo en entorno empresarial con MV3) \\
\hline
\rowcolor{suspiciousYellow}
\textbf{Sospechosos} & \texttt{tabs} (URLs), \texttt{history} (limpieza), \texttt{scripting}, \texttt{nativeMessaging} (VPNs) \\
\hline
\rowcolor{incoherentRed}
\textbf{Incoherentes} & \texttt{debugger}, \texttt{proxy} (excepto VPNs), \texttt{desktopCapture}, \texttt{tabCapture}, \texttt{pageCapture}, \texttt{identity} \\
\hline
\end{tabular}
\end{table}

\section{Dependencia Funcional de Permisos: Capacidad vs. Alcance}
\label{sec:dependencia_funcional}

Un hallazgo crítico del análisis de permisos es que varios permisos clasificados como de riesgo Alto no representan un riesgo efectivo por sí solos. Esto se debe a la distinción entre \textbf{capacidad} y \textbf{alcance} en el modelo de permisos de Chrome:

\begin{itemize}
    \item \textbf{Permisos de capacidad}: Habilitan una funcionalidad específica de la API, pero no otorgan acceso a ningún dominio. Ejemplos: \texttt{scripting} (habilita \texttt{chrome.scripting.\allowbreak executeScript()}), \texttt{cookies} (habilita \texttt{chrome.cookies.\allowbreak get/set()}), \texttt{declarativeNetRequest} (habilita \texttt{update\-Dynamic\-Rules()}).
    \item \textbf{Permisos de alcance}: Determinan sobre qué dominios puede operar la extensión. Incluyen \texttt{host\_permissions} (acceso permanente a dominios específicos o a todos) y \texttt{activeTab} (acceso temporal a la pestaña activa, condicionado a interacción del usuario).
\end{itemize}

El riesgo efectivo de un permiso de capacidad es función del alcance que lo acompaña:

\begin{equation}
\text{Riesgo}_{\text{efectivo}}(\text{capacidad}) = \text{Riesgo}_{\text{intrínseco}}(\text{capacidad}) \times f(\text{alcance})
\label{eq:riesgo_efectivo}
\end{equation}

Donde $f(\text{alcance}) = 0$ cuando la extensión no posee permisos de host, y $f(\text{alcance})$ escala con la amplitud de los dominios concedidos. En la fórmula de cálculo (Sección \ref{sec:formula}), esta función continua se implementa como la función discreta $f(H)$ definida en la Tabla \ref{tab:host_factor}, con cinco niveles que corresponden a los tipos de \texttt{host\_permissions} observados en extensiones reales. Esto implica que:

\begin{itemize}
    \item \texttt{scripting} sin \texttt{host\_permissions} ni \texttt{activeTab} $\Rightarrow$ riesgo efectivo $= 0$ (la API no puede ejecutarse sobre ninguna página).
    \item \texttt{cookies} sin \texttt{host\_permissions} $\Rightarrow$ solo accede a cookies del dominio de la propia extensión (riesgo mínimo).
    \item \texttt{declarativeNetRequest} con acción \texttt{modifyHeaders} requiere \texttt{host\_permissions} o \texttt{declarativeNetRequestWithHostAccess} para los dominios objetivo.
\end{itemize}

En contraste, existen \textbf{permisos estáticos} cuyo riesgo no depende de los permisos de host: \texttt{debugger} (acceso al protocolo DevTools), \texttt{nativeMessaging} (comunicación con binarios del SO), \texttt{management} (gestión de otras extensiones), entre otros.

Esta distinción es fundamental para el modelo de riesgo propuesto y se refleja directamente en la fórmula de cálculo (Sección \ref{sec:formula}), donde los permisos sensibles a host se ponderan multiplicativamente por un factor de alcance, mientras que los permisos estáticos mantienen su peso independientemente.

\textbf{Nota sobre \texttt{activeTab}:} Este permiso merece tratamiento especial. Aunque otorga acceso de host, lo hace de forma \textit{temporal} y \textit{condicional}: solo se activa tras una acción explícita del usuario (clic en el icono de la extensión) y se revoca al cambiar de pestaña o navegar. En el modelo propuesto, \texttt{activeTab} se trata como ``Host Temporal/Condicional'', con un factor de alcance intermedio entre no tener hosts y poseer \texttt{<all\_urls>}.

\section{Fórmula de Cálculo de Puntaje de Riesgo}
\label{sec:formula}

\subsection{Fundamento Metodológico}

El modelo de evaluación de riesgo propuesto adapta metodologías reconocidas internacionalmente al contexto específico de extensiones de navegador. Se fundamenta en dos estándares principales:

\textbf{NIST SP 800-30 Rev. 1 (2012)} \cite{nist_sp800_30}: La guía del Instituto Nacional de Estándares y Tecnología de Estados Unidos para conducir evaluaciones de riesgo en sistemas de información. Define el riesgo como una función de la probabilidad de que una amenaza explote una vulnerabilidad y el impacto adverso resultante:

\begin{quote}
``El riesgo es una función de la probabilidad de ocurrencia de un evento de amenaza y del impacto adverso potencial en caso de que el evento ocurra.''\footnote{Traducción propia. Original en inglés: ``Risk is a function of the likelihood of a threat event's occurrence and potential adverse impact should the event occur.''} \cite{nist_sp800_30}
\end{quote}

NIST SP 800-30 establece que las organizaciones pueden emplear enfoques cualitativos, cuantitativos o semi-cuantitativos, y que los factores de riesgo típicos incluyen: amenaza, vulnerabilidad, impacto, probabilidad y condiciones predisponentes.

\textbf{ISO/IEC 27005:2022} \cite{iso_27005}: El estándar internacional para gestión de riesgos de seguridad de la información. La versión 2022 formalizó explícitamente el análisis semi-cuantitativo como técnica válida dentro del estándar, donde algunos aspectos (como probabilidad) se cuantifican mediante métodos estadísticos y otros (como impacto) se definen mediante métodos subjetivos como opinión de expertos.

ISO 27005 soporta tanto análisis cualitativo (usando escalas descriptivas como alto/medio/bajo) como cuantitativo (usando valores numéricos como expectativa de pérdida).

\subsection{Adaptación al Contexto de Extensiones}

La fórmula propuesta adapta el modelo general de riesgo ($\text{Riesgo} = f(\text{Probabilidad}, \text{Impacto})$) al contexto de extensiones de navegador, incorporando la distinción entre capacidad y alcance descrita en la Sección \ref{sec:dependencia_funcional}. La estructura resultante es análoga a modelos de scoring reconocidos en ciberseguridad: CVSS (\textit{Common Vulnerability Scoring System}), que utiliza el producto de métricas de impacto y explotabilidad; CWSS (\textit{Common Weakness Scoring System}), cuya fórmula $\text{Score} = \text{Base} \times \text{Attack Surface} \times \text{Environmental}$ presenta una estructura multiplicativa similar a Peso $\times$ Coherencia $\times$ Alcance; y FAIR (\textit{Factor Analysis of Information Risk}), que descompone el riesgo en $\text{Loss Event Frequency} \times \text{Loss Magnitude}$. A diferencia de estos modelos genéricos, la fórmula propuesta incorpora la distinción específica entre permisos sensibles a host y estáticos, lo que permite capturar la arquitectura de seguridad particular de extensiones de navegador bajo Manifest V3:

\begin{equation}
\text{Riesgo} = \underbrace{\sum_{i \in S} (\text{Peso}_i \times \text{Factor}_{\text{Coherencia}_i} \times f(H))}_{\text{Permisos sensibles a host}} + \underbrace{\sum_{j \in E} (\text{Peso}_j \times \text{Factor}_{\text{Coherencia}_j})}_{\text{Permisos estáticos}}
\end{equation}

El modelo distingue dos tipos de permisos:

\begin{itemize}
    \item \textbf{Permisos sensibles a host} ($S$): Permisos cuyo riesgo efectivo depende del alcance de \texttt{host\_permissions}. Su peso se multiplica por $f(H)$, el factor de alcance.
    \item \textbf{Permisos estáticos} ($E$): Permisos cuyo riesgo es independiente de los hosts concedidos. Estos mantienen su peso íntegro sin factor de alcance.
\end{itemize}

La Tabla \ref{tab:se_classification} presenta la clasificación exhaustiva de cada permiso en uno de los dos conjuntos:

\begin{table}[!htbp]
\centering
\small
\begin{tabular}{|p{3cm}|p{5.5cm}|p{5.5cm}|}
\hline
\textbf{Nivel} & \textbf{Sensibles a host ($S$)} & \textbf{Estáticos ($E$)} \\
\hline
CRÍTICO & \texttt{tabCapture}, \texttt{pageCapture} & \texttt{debugger}, \texttt{nativeMessaging}, \texttt{proxy}, \texttt{vpnProvider} \\
\hline
ALTO & \texttt{cookies}, \texttt{scripting}, \texttt{declarativeNetRequest}, \texttt{webRequest}, \texttt{webRequestBlocking}, \texttt{userScripts}, \texttt{declarativeNetRequest\-WithHostAccess}, \texttt{desktopCapture} & \texttt{history}, \texttt{downloads}, \texttt{downloads.open}, \texttt{privacy}, \texttt{browsingData}, \texttt{contentSettings}, \texttt{webNavigation}, \texttt{webAuthentication\-Proxy}, \texttt{certificateProvider}, \texttt{platformKeys} \\
\hline
MEDIO & \texttt{activeTab} (host temporal) & \texttt{alarms}, \texttt{bookmarks}, \texttt{clipboard\-Read}, \texttt{clipboardWrite}, \texttt{geolocation}, \texttt{identity}, \texttt{management}, \texttt{sessions}, \texttt{topSites}, \texttt{context\-Menus}, \texttt{tabGroups}, \texttt{dns}, \texttt{tabs}, \texttt{offscreen}, \texttt{processes} \\
\hline
BAJO & --- & Todos los permisos de nivel Bajo (\texttt{storage}, \texttt{notifications}, \texttt{idle}, \texttt{power}, \texttt{tts}, etc.) \\
\hline
\end{tabular}
\caption{Clasificación exhaustiva de permisos: sensibles a host ($S$) vs. estáticos ($E$)}
\label{tab:se_classification}
\end{table}

Donde:
\begin{itemize}
    \item \textbf{Peso del Permiso} ($\text{Peso}_i$): Severidad del daño potencial si el permiso es abusado. Se deriva de la clasificación de Google (2019) actualizada con evidencia de ataques recientes y análisis cualitativos de riesgo de fuentes como LayerX (2025) y Singh et al. (2025).
    \item \textbf{Factor de Coherencia} ($\text{Factor}_{\text{Coherencia}_i}$): Probabilidad de que el permiso sea usado maliciosamente, estimada mediante la coherencia entre el permiso y la funcionalidad declarada de la extensión. \textbf{Esta es una contribución original} del proyecto: las matrices de correlación permiso-categoría no existían previamente.
    \item \textbf{Factor de Alcance} ($f(H)$): Escala el riesgo de permisos sensibles según la amplitud de los \texttt{host\_permissions} declarados (ver Tabla \ref{tab:host_factor}).
\end{itemize}

\subsection{Escalas de Evaluación}

Los pesos de impacto se asignan en escala semi-logarítmica, siguiendo las recomendaciones de NIST SP 800-30 para evaluaciones semi-cuantitativas que requieren priorización relativa. La escala utiliza una progresión donde cada nivel representa aproximadamente el doble del anterior (ratio $\approx 2\text{--}2.5\times$), similar a la empleada en modelos de scoring reconocidos como CWSS (\textit{Common Weakness Scoring System}):

\begin{table}[!htbp]
\centering
\begin{tabular}{|l|c|c|p{6cm}|}
\hline
\textbf{Nivel de Impacto} & \textbf{Peso} & \textbf{Ratio} & \textbf{Justificación} \\
\hline
CRÍTICO & 10 & --- & Permisos que permiten escape del sandbox o acceso irrestricto (Google \textit{Highest Risk}). \\
ALTO & 5 & $2\times$ & Permisos con acceso significativo a datos sensibles (Google \textit{High Risk} + elevaciones documentadas). \\
MEDIO & 2 & $2.5\times$ & Permisos que requieren evaluación contextual (Google \textit{Medium Risk}). \\
BAJO & 1 & $2\times$ & Permisos con impacto mínimo (Google \textit{Low Risk}). Cada permiso contribuye al menos 1 unidad base al puntaje. \\
\hline
\end{tabular}
\caption{Pesos de impacto basados en clasificación de Google actualizada}
\label{tab:impact_weights}
\end{table}

Los factores de coherencia (probabilidad de abuso) se basan en las matrices de correlación permiso-categoría desarrolladas en este proyecto. La escala utiliza el nivel ``Sospechoso'' como línea base ($1.0$), con atenuación para permisos esperados y amplificación creciente para permisos incoherentes:

\begin{table}[!htbp]
\centering
\begin{tabular}{|l|c|c|p{5.5cm}|}
\hline
\textbf{Coherencia} & \textbf{Factor} & \textbf{Ratio} & \textbf{Justificación} \\
\hline
Esperado & 0.1 & $\div 10\times$ & Permiso coherente con la funcionalidad de la categoría. Probabilidad de abuso muy baja (uso legítimo esperado). \\
Sospechoso & 1.0 & base & Permiso que requiere justificación adicional. Probabilidad de abuso moderada (podría ser legítimo o no). \\
Incoherente & 5.0 & $\times 5$ & Permiso sin justificación razonable para la categoría. Probabilidad de abuso alta (indicador de posible intención maliciosa). \\
\hline
\end{tabular}
\caption{Factores de coherencia (probabilidad de abuso)}
\label{tab:coherence_factors}
\end{table}

El factor de alcance $f(H)$ refleja el principio de mínimo privilegio y escala el riesgo de permisos sensibles a host según la amplitud de los dominios concedidos:

\begin{table}[!htbp]
\centering
\begin{tabular}{|l|c|p{7cm}|}
\hline
\textbf{Tipo de Host} & $f(H)$ & \textbf{Justificación} \\
\hline
Sin \texttt{host\_permissions} & 0.0 & El permiso de capacidad es inerte; no puede operar sobre ningún dominio externo. \\
Solo \texttt{activeTab} & 0.3 & Host temporal y condicional: requiere interacción del usuario, se revoca al navegar. Riesgo limitado por diseño. \\
Dominios específicos & 0.5 & Alcance acotado a sitios declarados. Riesgo proporcional a la sensibilidad de los dominios. \\
Wildcards amplios (\texttt{*://*/*}) & 0.8 & Acceso a todos los sitios de un protocolo. Superficie de ataque extensa. \\
\texttt{<all\_urls>} o equivalente & 1.0 & Acceso universal. Máxima superficie de ataque. \\
\hline
\end{tabular}
\caption{Factor de alcance $f(H)$ para permisos sensibles a host}
\label{tab:host_factor}
\end{table}

Esta estructura multiplicativa captura un principio fundamental: \textbf{el riesgo real de una extensión no es la suma de sus capacidades teóricas, sino la combinación de lo que puede hacer (capacidad) con dónde puede hacerlo (alcance)}. Una extensión con \texttt{scripting} + \texttt{cookies} pero sin \texttt{host\_permissions} tiene riesgo efectivo $\approx 0$ para esos permisos, mientras que la misma combinación con \texttt{<all\_urls>} representa un riesgo máximo.

\subsection{Calibración Planificada}

Los valores numéricos propuestos (pesos 10/5/2/1, factores de coherencia 0.1/1.0/5.0, y factores de alcance 0.0--1.0) son valores iniciales con progresión lógica documentada que serán calibrados en la Fase 3 del proyecto mediante análisis de extensiones reales del Chrome Web Store. El proceso de calibración incluirá:

\begin{enumerate}
    \item Recolección de datos de al menos 1,000 extensiones con etiquetas conocidas (legítimas vs. removidas por malware).
    \item Ajuste de pesos mediante optimización supervisada, utilizando la separación de clases (legítimas vs. maliciosas) como función objetivo. Se empleará análisis ROC (\textit{Receiver Operating Characteristic}) para determinar los umbrales óptimos de clasificación, y regresión logística para validar la capacidad discriminante de los puntajes resultantes.
    \item Validación cruzada (\textit{k-fold}) para evitar sobreajuste.
\end{enumerate}

\subsection{Interpretación del Puntaje}

\begin{table}[!htbp]
\centering
\begin{tabular}{|c|l|p{8cm}|}
\hline
\textbf{Rango} & \textbf{Nivel} & \textbf{Interpretación} \\
\hline
0--10 & Bajo & Extensión con permisos coherentes y limitados. \\
\hline
11--25 & Moderado & Requiere revisión de justificación de permisos. \\
\hline
26--50 & Alto & Extensión con permisos que exceden su funcionalidad aparente. \\
\hline
51+ & Crítico & Extensión potencialmente maliciosa o con permisos incoherentes. \\
\hline
\end{tabular}
\caption{Interpretación de puntajes de riesgo}
\end{table}

\subsection{Ejemplos de Cálculo}

Para validar que los umbrales de interpretación producen clasificaciones coherentes, se presentan tres ejemplos de cálculo representativos.

\textbf{Ejemplo 1: Extensión legítima --- Ad-blocker (categoría Privacy \& Security)}

Permisos típicos: \texttt{storage} (Bajo, Esperado), \texttt{activeTab} (Medio, Esperado), \texttt{declarativeNetRequest} (Alto, Esperado). Host: solo \texttt{activeTab} $\Rightarrow f(H)=0.3$.

\begin{itemize}
    \item \texttt{storage} (estático): $1 \times 0.1 = 0.1$
    \item \texttt{activeTab} (estático): $2 \times 0.1 = 0.2$
    \item \texttt{declarativeNetRequest} (sensible a host): $5 \times 0.1 \times 0.3 = 0.15$
\end{itemize}
\textbf{Puntaje total: 0.45} $\Rightarrow$ \textbf{Bajo.} Resultado coherente: un ad-blocker legítimo con permisos esperados tiene riesgo mínimo.

\textbf{Ejemplo 2: Extensión sospechosa --- Herramienta ``productividad'' con permisos excesivos (categoría Tools)}

Permisos: \texttt{storage} (Bajo, Esperado), \texttt{cookies} (Alto, Incoherente), \texttt{scripting} (Alto, Sospechoso), \texttt{history} (Alto, Incoherente), \texttt{alarms} (Medio, Esperado). Host: \texttt{<all\_urls>} $\Rightarrow f(H)=1.0$.

\begin{itemize}
    \item \texttt{storage} (estático): $1 \times 0.1 = 0.1$
    \item \texttt{cookies} (sensible a host): $5 \times 5.0 \times 1.0 = 25.0$
    \item \texttt{scripting} (sensible a host): $5 \times 1.0 \times 1.0 = 5.0$
    \item \texttt{history} (estático): $5 \times 5.0 = 25.0$
    \item \texttt{alarms} (estático): $2 \times 0.1 = 0.2$
\end{itemize}
\textbf{Puntaje total: 55.3} $\Rightarrow$ \textbf{Crítico.} Resultado coherente: una extensión de ``herramientas'' que solicita cookies y history con acceso a todos los sitios es altamente sospechosa.

\textbf{Ejemplo 3: Extensión tipo DataByCloud (categoría Tools, maliciosa conocida)}

Permisos documentados: \texttt{cookies} (Alto, Incoherente), \texttt{management} (Medio, Incoherente), \texttt{scripting} (Alto, Incoherente), \texttt{storage} (Bajo, Esperado), \texttt{declarativeNetRequest} (Alto, Incoherente). Host: dominios específicos (Workday, NetSuite, SuccessFactors) $\Rightarrow f(H)=0.5$.

\begin{itemize}
    \item \texttt{cookies} (sensible a host): $5 \times 5.0 \times 0.5 = 12.5$
    \item \texttt{management} (estático): $2 \times 5.0 = 10.0$
    \item \texttt{scripting} (sensible a host): $5 \times 5.0 \times 0.5 = 12.5$
    \item \texttt{storage} (estático): $1 \times 0.1 = 0.1$
    \item \texttt{declarativeNetRequest} (sensible a host): $5 \times 5.0 \times 0.5 = 12.5$
\end{itemize}
\textbf{Puntaje total: 47.6} $\Rightarrow$ \textbf{Alto} (cercano a Crítico). Resultado coherente: aun con hosts limitados (no \texttt{<all\_urls>}), la combinación de permisos incoherentes para una extensión de herramientas genera una alerta significativa. Con \texttt{<all\_urls>} ($f(H)=1.0$), el puntaje ascendería a 85.1 (Crítico).

Estos ejemplos demuestran que los umbrales propuestos separan adecuadamente extensiones legítimas (Bajo) de extensiones con patrones de permisos maliciosos conocidos (Alto/Crítico).

\endgroup


% ============================================
% RESULTADOS ESPERADOS
% ============================================

\chapter{Resultados Esperados}
\label{chap:resultados}

\begingroup
\setlength{\parindent}{0pt}
\setlength{\parskip}{0.6em}

Una extensión de seguridad completamente funcional, compatible con la arquitectura \textit{Manifest V3} de Google Chrome. Si bien la arquitectura de extensiones es compartida por otros navegadores basados en Chromium, la validación se realizará exclusivamente en Google Chrome. Este artefacto de software integrará:
\begin{itemize}
    \item Un Motor de Auditoría capaz de asignar un puntaje de riesgo a las extensiones instaladas basándose en la coherencia entre sus permisos y su categoría funcional.
    \item Un módulo de Defensa contra Actualizaciones (Update Watcher) que alerte al usuario sobre cambios de privilegios o comportamiento tras una actualización automática.
    \item Un sistema de Control Contextual que permita al usuario definir "Zonas Seguras" (ej. Banca, Trabajo) donde se restrinja dinámicamente la ejecución de otras extensiones.
\end{itemize}

\textbf{Nota sobre permisos de la propia herramienta:} La extensión propuesta requiere permisos como \texttt{management} (nivel Medio), \texttt{tabs} y \texttt{storage} para operar. El modelo de riesgo clasificaría estos permisos como coherentes con la categoría \textit{Privacy \& Security}, según las matrices de correlación definidas en el capítulo \ref{chap:modelo_riesgo}. La herramienta se excluirá automáticamente de su propio análisis para evitar alertas recursivas, y su puntaje de riesgo será transparente para el usuario desde la interfaz de configuración.

\textbf{Limitación conocida --- categoría autodeclarada:} El modelo de riesgo asume que la categoría declarada por la extensión en la Chrome Web Store refleja su funcionalidad real. Sin embargo, esta categoría es elegida libremente por el desarrollador y podría ser manipulada. Una extensión maliciosa podría registrarse en la categoría \textit{Privacy \& Security} para que permisos como \texttt{cookies} y \texttt{webRequest} se clasifiquen como ``esperados'', reduciendo artificialmente su puntaje de riesgo. La validación cruzada de la categoría declarada con la funcionalidad observada (por ejemplo, analizando las APIs efectivamente invocadas) se identifica como una mejora futura que fortalecería la robustez del modelo.

\endgroup
% =========================
% BIBLIOGRAFÍA
% =========================
\begin{thebibliography}{99}

\bibitem{thn2025shadypanda}
The Hacker News, 
``A Browser Extension Risk Guide After the ShadyPanda Campaign,'' 
Dec. 2025. 
[Online]. Available: \url{https://thehackernews.com/2025/12/a-browser-extension-risk-guide-after.html}

\bibitem{pandya2026databycloud}
Socket Threat Research Team (K. Pandya), 
``5 Malicious Chrome Extensions Enable Session Hijacking in Enterprise HR and ERP Systems,'' 
Socket, Jan. 2026. 
[Online]. Available: \url{https://socket.dev/blog/5-malicious-chrome-extensions-enable-session-hijacking}

\bibitem{singh2025study}
S. Singh, G. Varshney, T. K. Singh, V. Mishra, and K. Verma, 
``A Study on Malicious Browser Extensions in 2025,'' 
arXiv:2503.04292v2 [cs.CR], Mar. 2025 (rev. Oct. 2025).
Department of CSE, IIT Jammu, India.
[Online]. Available: \url{https://arxiv.org/abs/2503.04292}

\bibitem{dardikman2025aiconversations}
I. Dardikman, 
``8 Million Users' AI Conversations Sold for Profit by 'Privacy' Extensions,'' 
Koi Research, Dec. 2025. 
[Online]. Available: \url{https://www.koi.ai/blog/urban-vpn-browser-extension-ai-conversations-data-collection}

\bibitem{zenitylabs2025claude}
R. Klugman-Onitza and J. Donato, 
``Claude in Chrome: A Threat Analysis,'' 
Zenity Labs, Dec. 2025. 
[Online]. Available: \url{https://labs.zenity.io/p/claude-in-chrome-a-threat-analysis}

\bibitem{chrome_extensions_overview}
Google, 
``Extensions architecture overview,'' 
\textit{Chrome for Developers}, 
[Online]. Available: \url{https://developer.chrome.com/docs/extensions/mv3/architecture-overview/}

\bibitem{chrome_extensions_security}
Google, 
``Stay secure,'' 
\textit{Chrome for Developers}, 
[Online]. Available: \url{https://developer.chrome.com/docs/extensions/mv3/security/}

\bibitem{isolated_worlds}
Google, 
``Content scripts: Isolated worlds,'' 
\textit{Chrome for Developers}, 
[Online]. Available: 
\url{https://developer.chrome.com/docs/extensions/mv3/content_scripts/#isolated_worlds}

\bibitem{pantelaios2024manifestv3}
N. Pantelaios and A. Kapravelos, 
``Work-in-Progress: Manifest V3 Unveiled: Navigating the New Era of Browser Extensions,'' 
in \textit{NDSS Symposium}, 2024.

\bibitem{manifestv3chrome}
Google, 
``Manifest V3,'' 
\textit{Chrome for Developers}, 
[Online]. Available: \url{https://developer.chrome.com/docs/extensions/mv3/intro/}

\bibitem{declarative_net_request}
Google, 
``Modify network requests with declarativeNetRequest,'' 
\textit{Chrome for Developers}, 
[Online]. Available: \url{https://developer.chrome.com/docs/extensions/reference/declarativeNetRequest/}

\bibitem{service_workers}
Google, 
``About extension service workers,'' 
\textit{Chrome for Developers}, 
[Online]. Available: \url{https://developer.chrome.com/docs/extensions/mv3/service_workers/}

\bibitem{koi2025darkspectre}
Koi Research, 
``DarkSpectre: Unmasking the Threat Actor Behind 8.8 Million Infected Browsers,'' 
Tech. Rep., 2025.

\bibitem{seetharam2025genai}
S. B. Seetharam, M. Nabeel, and W. Melicher,
``Malicious GenAI Chrome Extensions: Unpacking Data Exfiltration and Malicious Behaviours,''
Palo Alto Networks, presented at \textit{Virus Bulletin Conference}, Dec. 2025.
arXiv:2512.10029 [cs.CR].
[Online]. Available: \url{https://arxiv.org/abs/2512.10029}

\bibitem{stealth_extension}
C. Lim, B. Jin, and H. Kim,
``Stealth Extension Exfiltration (SEE) Attacks: Stealing User Data without Permissions via Browser Extensions,''
in \textit{Proc. 40th ACM/SIGAPP Symposium on Applied Computing (SAC '25)},
Mar.--Apr. 2025, pp. 1--8.
[Online]. Available: \url{https://doi.org/10.1145/3672608.3707856}

\bibitem{yu2023coco}
J. Yu, S. Li, J. Zhu, and Y. Cao, 
``CoCo: Efficient Browser Extension Vulnerability Detection via Coverage-guided, 
Concurrent Abstract Interpretation,'' 
in \textit{Proc. 2023 ACM SIGSAC Conf. Computer and Communications Security (CCS '23)}, 
Nov. 2023, pp. 2441--2455.

\bibitem{doublex}
A. Fass, D. F. Somé, M. Backes, and B. Stock,
``DoubleX: Statically Detecting Vulnerable Data Flows in Browser Extensions at Scale,''
in \textit{Proc. ACM SIGSAC Conf. Computer and Communications Security (CCS)},
Nov. 2021, pp. 1789--1804.
[Online]. Available: \url{https://doi.org/10.1145/3460120.3484745}

\bibitem{sam2023mitigating}
J. Sam and J. A. Jenifer,
``Mitigating the Security Risks of Browser Extensions,''
in \textit{2023 Int. Conf. Sustainable Computing and Smart Systems (ICSCSS)},
Jun. 2023, pp. 1460--1465.
doi: 10.1109/ICSCSS57650.2023.10169483.

\bibitem{arcanum}
Q. Xie, M. V. Kasi Murali, P. Pearce, and F. Li,
``Arcanum: Detecting and Evaluating the Privacy Risks of Browser Extensions on Web Pages and Web Content,''
in \textit{33rd USENIX Security Symposium (USENIX Security 24)},
Aug. 2024, pp. 4607--4624.
[Online]. Available: \url{https://www.usenix.org/conference/usenixsecurity24/presentation/xie-qinge}

\bibitem{kathir2026vkstyles}
A. Cohen, 
``VK Styles: 500K Users Infected by Chrome Extensions That Hijack VKontakte Accounts,'' 
Koi Security, Feb. 2026.
[Online]. Available: \url{https://www.koi.ai/blog/vk-styles-500k-users-infected-by-chrome-extensions-that-hijack-vkontakte-accounts}

\bibitem{pantelaios2020malicious}
N. Pantelaios, N. Nikiforakis, and A. Kapravelos, 
``You've Changed: Detecting Malicious Browser Extensions through Their Update Deltas,'' 
in \textit{Proc. ACM SIGSAC Conf. Computer and Communications Security (CCS)}, 
2020, pp. 477--491.

% Referencias para el capítulo de Modelo de Riesgo

\bibitem{google_permission_risk}
Google Chrome Enterprise \& Education,
``Permission Risk Whitepaper: Understand the risks of permissions for Chrome extensions,''
Jul. 2019.
[Online]. Available: \url{https://support.google.com/chrome/a/answer/9897812}

\bibitem{nayak2024experimental}
A. Nayak, R. Khandelwal, E. Fernandes, and K. Fawaz,
``Experimental Security Analysis of Sensitive Data Access by Browser Extensions,''
in \textit{Proc. ACM Web Conference 2024 (WWW '24)}, May 2024, pp. 1--12.

\bibitem{layerx2025report}
LayerX Security,
``Enterprise Browser Extension Security Report 2025,''
Tech. Rep., Apr. 2025.
[Online]. Available: \url{https://go.layerxsecurity.com/hubfs/LayerX_Enterprise_Browser_Extension_Security_Report_2025.pdf}

\bibitem{chrome_alarms_docs}
Google,
``chrome.alarms API,''
\textit{Chrome for Developers},
[Online]. Available: \url{https://developer.chrome.com/docs/extensions/reference/api/alarms}

\bibitem{socket2025phantom}
Socket Research,
``Two Chrome Extensions Caught Secretly Stealing Credentials from Over 170 Sites,''
Dec. 2025.
[Online]. Available: \url{https://socket.dev/blog/phantom-shuttle-chrome-extension-malware}

% Referencias para Marco Regulatorio

\bibitem{cws_policies}
Google,
``Chrome Web Store Program Policies,''
\textit{Chrome for Developers}, May 2025.
[Online]. Available: \url{https://developer.chrome.com/docs/webstore/program-policies}

\bibitem{eu_dsa}
European Commission,
``Digital Services Act: Questions and Answers,''
\textit{Shaping Europe's Digital Future}, 2024.
[Online]. Available: \url{https://digital-strategy.ec.europa.eu/en/faqs/digital-services-act-questions-and-answers}

\bibitem{eu_digital_omnibus}
European Commission,
``Digital Omnibus Package: Simplifying EU Digital Rules,''
Nov. 2025.
[Online]. Available: \url{https://digital-strategy.ec.europa.eu/en/faqs/digital-package}

\bibitem{chrome_permissions_list}
Google,
``Permissions List,''
\textit{Chrome for Developers}, Apr. 2025.
[Online]. Available: \url{https://developer.chrome.com/docs/extensions/reference/permissions-list}

% Referencias para justificación metodológica

\bibitem{nist_sp800_30}
National Institute of Standards and Technology,
``Guide for Conducting Risk Assessments,''
NIST Special Publication 800-30 Revision 1, Sep. 2012.
[Online]. Available: \url{https://csrc.nist.gov/pubs/sp/800/30/r1/final}

\bibitem{iso_27005}
International Organization for Standardization,
``ISO/IEC 27005:2022 Information security, cybersecurity and privacy protection --- Guidance on managing information security risks,''
4th ed., 2022.

\end{thebibliography}


\end{document}